\documentclass[12pt, letterpaper, oneside]{book}

% ──────────────────────────────────────────────────────────────
% PACKAGES
% ──────────────────────────────────────────────────────────────
\usepackage[margin=1in]{geometry}
\usepackage[utf8]{inputenc}
\usepackage[T1]{fontenc}
\usepackage{lmodern}
\usepackage{microtype}
\usepackage{graphicx}
\usepackage{xcolor}
\usepackage{booktabs}
\usepackage{longtable}
\usepackage{tabularx}
\usepackage{array}
\usepackage{enumitem}
\usepackage{tikz}
\usepackage{fancyhdr}
\usepackage{titlesec}
\usepackage{tocloft}
\usepackage{hyperref}
\usepackage{cleveref}
\usepackage{amssymb}
\usepackage{pifont}
\usepackage{framed}
\usepackage{mdframed}
\usepackage{float}
\usepackage{caption}
\usepackage{multicol}
\usepackage{calc}
\usepackage{etoolbox}

% ──────────────────────────────────────────────────────────────
% COLORS
% ──────────────────────────────────────────────────────────────
\definecolor{primary}{HTML}{7E22CE}
\definecolor{accent}{HTML}{F97316}
\definecolor{darkbg}{HTML}{1E293B}
\definecolor{success}{HTML}{16A34A}
\definecolor{danger}{HTML}{DC2626}
\definecolor{info}{HTML}{2563EB}
\definecolor{lightgray}{HTML}{F1F5F9}
\definecolor{midgray}{HTML}{94A3B8}
\definecolor{money}{HTML}{16A34A}

% ──────────────────────────────────────────────────────────────
% HYPERREF SETUP
% ──────────────────────────────────────────────────────────────
\hypersetup{
  colorlinks=true,
  linkcolor=primary,
  urlcolor=info,
  citecolor=primary,
  bookmarks=true,
  bookmarksopen=true,
  pdftitle={Crayons \& Quills: Complete Business Launch Guide},
  pdfauthor={Crayons \& Quills --- an Atman Labs LLC product},
  pdfsubject={AI-Powered Children's Book Business},
}

% ──────────────────────────────────────────────────────────────
% CUSTOM ENVIRONMENTS
% ──────────────────────────────────────────────────────────────
\newmdenv[
  backgroundcolor=lightgray,
  linecolor=primary,
  linewidth=2pt,
  topline=false,
  bottomline=false,
  rightline=false,
  innerleftmargin=15pt,
  innerrightmargin=15pt,
  innertopmargin=10pt,
  innerbottommargin=10pt,
  skipabove=12pt,
  skipbelow=12pt
]{infobox}

\newmdenv[
  backgroundcolor=red!5,
  linecolor=danger,
  linewidth=2pt,
  topline=false,
  bottomline=false,
  rightline=false,
  innerleftmargin=15pt,
  innerrightmargin=15pt,
  innertopmargin=10pt,
  innerbottommargin=10pt,
  skipabove=12pt,
  skipbelow=12pt
]{warningbox}

\newmdenv[
  backgroundcolor=green!5,
  linecolor=success,
  linewidth=2pt,
  topline=false,
  bottomline=false,
  rightline=false,
  innerleftmargin=15pt,
  innerrightmargin=15pt,
  innertopmargin=10pt,
  innerbottommargin=10pt,
  skipabove=12pt,
  skipbelow=12pt
]{successbox}

\newmdenv[
  backgroundcolor=orange!5,
  linecolor=accent,
  linewidth=2pt,
  topline=false,
  bottomline=false,
  rightline=false,
  innerleftmargin=15pt,
  innerrightmargin=15pt,
  innertopmargin=10pt,
  innerbottommargin=10pt,
  skipabove=12pt,
  skipbelow=12pt
]{costbox}

\newmdenv[
  backgroundcolor=blue!3,
  linecolor=info,
  linewidth=2pt,
  topline=false,
  bottomline=false,
  rightline=false,
  innerleftmargin=15pt,
  innerrightmargin=15pt,
  innertopmargin=10pt,
  innerbottommargin=10pt,
  skipabove=12pt,
  skipbelow=12pt
]{actionbox}

% ──────────────────────────────────────────────────────────────
% CUSTOM COMMANDS
% ──────────────────────────────────────────────────────────────
\newcommand{\cost}[1]{\textcolor{money}{\textbf{#1}}}
\newcommand{\cmark}{\ding{51}}
\newcommand{\xmark}{\ding{55}}
\newcommand{\step}[1]{\noindent\textbf{\textcolor{primary}{Step #1:}}~}
\newcommand{\substep}[1]{\noindent\hspace{1em}\textbf{\textcolor{midgray}{#1.}}~}

% ──────────────────────────────────────────────────────────────
% PAGE STYLE
% ──────────────────────────────────────────────────────────────
\pagestyle{fancy}
\fancyhf{}
\fancyhead[L]{\small\textcolor{midgray}{\leftmark}}
\fancyhead[R]{\small\textcolor{midgray}{Crayons \& Quills Launch Guide}}
\fancyfoot[C]{\thepage}
\renewcommand{\headrulewidth}{0.4pt}
\renewcommand{\headrule}{\hbox to\headwidth{\color{midgray}\leaders\hrule height \headrulewidth\hfill}}

% ──────────────────────────────────────────────────────────────
% TITLE FORMATTING
% ──────────────────────────────────────────────────────────────
\titleformat{\chapter}[display]
  {\normalfont\huge\bfseries\color{primary}}
  {\chaptertitlename\ \thechapter}{20pt}{\Huge}
\titleformat{\section}
  {\normalfont\Large\bfseries\color{darkbg}}
  {\thesection}{1em}{}
\titleformat{\subsection}
  {\normalfont\large\bfseries\color{darkbg}}
  {\thesubsection}{1em}{}

% ──────────────────────────────────────────────────────────────
% BEGIN DOCUMENT
% ──────────────────────────────────────────────────────────────
\begin{document}

% ══════════════════════════════════════════════════════════════
% TITLE PAGE
% ══════════════════════════════════════════════════════════════
\begin{titlepage}
\begin{center}
\vspace*{1.5in}

{\fontsize{42}{50}\selectfont\bfseries\textcolor{primary}{Crayons \& Quills}}

\vspace{0.3in}
{\fontsize{18}{22}\selectfont\textcolor{darkbg}{Complete Business Launch Guide}}

\vspace{0.5in}
{\large\textcolor{midgray}{From Existing Codebase to a Live, AI-Agent-Operated\\
Personalized Children's Book Company}}

\vspace{0.2in}
{\normalsize\textcolor{midgray}{An Atman Labs LLC Product}}

\vspace{0.4in}
\rule{3in}{0.5pt}

\vspace{0.4in}
{\large\textcolor{darkbg}{A Step-by-Step Manual for Non-Technical Founders}}

\vspace{0.3in}
{\normalsize\textcolor{midgray}{Budget: \$2{,}000 \quad|\quad Location: South Carolina\\[6pt]
No Programming Experience Required}}

\vfill
{\small\textcolor{midgray}{Version 1.0 --- February 2026}}
\end{center}
\end{titlepage}

% ══════════════════════════════════════════════════════════════
% COPYRIGHT PAGE
% ══════════════════════════════════════════════════════════════
\thispagestyle{empty}
\vspace*{\fill}
\noindent\textcolor{midgray}{\small
This document is a comprehensive operational guide for launching and running\\
Crayons \& Quills (a product of Atman Labs LLC) as an automated business. All pricing data is based on publicly\\
available information as of February 2026 and is subject to change.\\[12pt]
Always verify current pricing before making purchasing decisions.\\[12pt]
This guide does not constitute legal, financial, or tax advice. Consult\\
with a licensed attorney and CPA for matters specific to your situation.}
\newpage

% ══════════════════════════════════════════════════════════════
% TABLE OF CONTENTS
% ══════════════════════════════════════════════════════════════
\tableofcontents
\newpage

% ══════════════════════════════════════════════════════════════
% EXECUTIVE SUMMARY
% ══════════════════════════════════════════════════════════════
\chapter*{Executive Summary}
\addcontentsline{toc}{chapter}{Executive Summary}

\section*{What You Have Already Built}

You already have a \textbf{fully architected codebase} for Crayons \& Quills. This is not a ``start from scratch'' guide---this is a ``configure, deploy, and launch'' guide. Your \texttt{LLM-books} repository contains:

\begin{itemize}[leftmargin=2em, itemsep=2pt]
  \item A \textbf{Next.js 14 frontend} with a book creation wizard, marketing landing pages, user dashboard, and responsive Tailwind CSS design
  \item A \textbf{NestJS backend} with JWT authentication, Prisma ORM, role-based access control, and versioned REST API endpoints
  \item A \textbf{complete AI services layer}: StoryEngineService, IllustrationEngineService, CharacterEngineService, and PdfAssemblyService
  \item \textbf{BullMQ background workers} for asynchronous book generation and illustration pipelines
  \item A \textbf{Prisma database schema} with 14 models (users, children, books, orders, payments, subscriptions, etc.)
  \item \textbf{Stripe payment integration}, \textbf{SendGrid email templates}, and \textbf{Lulu Direct fulfillment} provider
  \item \textbf{Docker Compose} configurations for both development and production
  \item An \textbf{agents/} directory with stub scripts for six Mac Mini automation agents
  \item \textbf{Claude Code configuration} files (\texttt{CLAUDE.md} and \texttt{.claude/settings.json}) for seamless development sessions
\end{itemize}

\noindent When fully configured and deployed, a customer visits your website, enters details about their child (name, age, interests, favorite animals, etc.), and your system \textbf{automatically}:

\begin{enumerate}[leftmargin=2em]
  \item Generates a unique story tailored to that child
  \item Creates custom illustrations for every other page
  \item Assembles a print-ready PDF
  \item Sends the file to a print-on-demand partner
  \item Ships a physical book to the customer's door
  \item Sends confirmation and tracking emails
\end{enumerate}

\noindent Your role is \textbf{business owner and supervisor}, not programmer or operator. AI agents handle the day-to-day: generating books, processing orders, sending emails, and monitoring for problems. You check a dashboard, handle edge cases, and focus on marketing.

\begin{successbox}
\textbf{Where you are today:} The codebase is built. What remains is configuring your third-party accounts (Chapters 1--4), setting your environment variables and deploying (Chapters 5--6), setting up automation (Chapters 7--8), connecting your print partner and payments (Chapters 9--10), and launching (Chapter 15). This guide walks you through each step from \textit{exactly where you are right now}.
\end{successbox}

\section*{What ``Run by AI Agents'' Actually Means}

Let us be clear about what is automated and what is not:

\begin{successbox}
\textbf{What AI Agents Handle Automatically:}
\begin{itemize}[leftmargin=1.5em, itemsep=2pt]
  \item Story and illustration generation (OpenAI API)
  \item PDF assembly and quality checks
  \item Order submission to print partner (Lulu)
  \item Payment processing (Stripe)
  \item Customer email notifications (order confirmation, shipping updates)
  \item Subscription renewals and reminders
  \item Fulfillment status tracking
\end{itemize}
\end{successbox}

\begin{warningbox}
\textbf{What Still Requires You (Estimated 2--5 Hours/Week):}
\begin{itemize}[leftmargin=1.5em, itemsep=2pt]
  \item Reviewing flagged orders or failed generations (rare)
  \item Handling customer support emails (consider AI chatbot later)
  \item Processing refund requests
  \item Monitoring your dashboard for errors
  \item Marketing: social media posts, ad campaigns, partnerships
  \item Accounting and taxes (quarterly, with a CPA)
  \item Updating pricing or adding new features over time
\end{itemize}
\end{warningbox}

\section*{Your \$2{,}000 Budget at a Glance}

\begin{table}[H]
\centering
\begin{tabularx}{\textwidth}{Xrr}
\toprule
\textbf{Category} & \textbf{One-Time} & \textbf{Monthly} \\
\midrule
South Carolina LLC formation & \$125 & \$0 \\
Domain name (1 year) & \$11 & --- \\
Hosting (Railway.app Pro) & --- & \$20 \\
AI coding tool (Claude Code via Claude Max) & --- & \$0 (already subscribed) \\
Automation platform (Make.com Core) & --- & \$11 \\
Email service (SendGrid Free) & --- & \$0 \\
Payment processing (Stripe) & --- & Per-transaction \\
OpenAI API (prepaid credits) & \$50 & Variable \\
Marketing budget (initial) & \$200 & Variable \\
Reserve for unexpected costs & \$300 & --- \\
\midrule
\textbf{Total Initial Outlay} & \textbf{\$686} & --- \\
\textbf{Monthly Operating Costs (pre-revenue)} & --- & \textbf{\$31} \\
\textbf{Remaining Reserve} & \textbf{\$1{,}314} & --- \\
\bottomrule
\end{tabularx}
\caption{Startup budget allocation summary}
\end{table}

\noindent At \$31/month in fixed costs, your \$1{,}314 reserve covers \textbf{42+ months} of operation even with zero revenue. You do not need to rush.

% ══════════════════════════════════════════════════════════════
% CHAPTER 1: FORM YOUR LLC
% ══════════════════════════════════════════════════════════════
\chapter{Form Your South Carolina LLC}

An LLC (Limited Liability Company) separates your personal assets from your business. If a customer ever sues over a book, they sue the LLC---not you personally. This chapter walks you through every step.

\section{Before You File: Choose Your Business Name}

\step{1} Pick 3--5 potential business names. Your LLC name does not have to match your website name, but it helps. Examples:
\begin{itemize}[leftmargin=2em]
  \item Atman Labs LLC
  \item Atman Labs Holdings LLC
  \item [YourCreativeName] LLC
\end{itemize}

\step{2} Search if the name is available:
\begin{actionbox}
Go to \url{https://businessfilings.sc.gov/BusinessFiling/Entity/Search}\\[4pt]
Type your desired name. If it shows \textbf{no results}, the name is available.\\[4pt]
Also check if a matching \texttt{.com} domain is available at \url{https://www.cloudflare.com/products/registrar/}
\end{actionbox}

\step{3} Optional: Reserve the name for 120 days by filing a Name Reservation with the SC Secretary of State (\cost{\$25} fee). This is optional---you can skip straight to filing.

\section{File Your Articles of Organization}

This is the official document that creates your LLC.

\step{4} Go to the SC Secretary of State online filing portal:
\begin{actionbox}
\url{https://businessfilings.sc.gov/BusinessFiling}\\[4pt]
Click ``Start a New Business'' $\rightarrow$ ``Limited Liability Company''
\end{actionbox}

\step{5} Fill in the required information:
\begin{itemize}[leftmargin=2em]
  \item \textbf{LLC Name:} Your chosen name (must end with ``LLC'' or ``L.L.C.'')
  \item \textbf{Registered Agent:} Enter \textbf{your own name} and your \textbf{South Carolina home address}. You are allowed to be your own registered agent as long as you are a SC resident over 18 with a physical street address (not a P.O.\ box). This saves \$50--125/year.
  \item \textbf{Organizer:} Your name and address
  \item \textbf{Management:} Choose ``Member-Managed'' (you manage it yourself)
  \item \textbf{Effective Date:} Today's date (or a future date if you prefer)
\end{itemize}

\step{6} Pay the filing fee: \cost{\$110} by mail or \cost{\$125} online (includes \$15 portal fee).

\begin{infobox}
\textbf{Processing time:} Online filings are typically processed within 3--5 business days. You will receive a confirmation with your LLC's file number.
\end{infobox}

\begin{costbox}
\textbf{South Carolina Advantage:} SC does \textbf{not} require annual reports or annual fees for standard LLCs taxed as sole proprietorships. Your total state cost after formation is \textbf{\$0/year}. This is one of the cheapest states to maintain an LLC.
\end{costbox}

\section{Get Your EIN (Employer Identification Number)}

An EIN is like a Social Security number for your business. You need it to open a business bank account, set up Stripe, and file taxes.

\step{7} Apply online at the IRS website:
\begin{actionbox}
\url{https://www.irs.gov/businesses/small-businesses-self-employed/apply-for-an-employer-identification-number-ein-online}\\[4pt]
Click ``Apply Online Now.'' Follow the prompts. Select:
\begin{itemize}[leftmargin=1.5em]
  \item Entity type: \textbf{Limited Liability Company}
  \item Number of members: \textbf{1}
  \item Reason: \textbf{Started new business}
  \item State: \textbf{South Carolina}
\end{itemize}
\end{actionbox}

\step{8} Print or save the confirmation letter with your EIN. You will need this number repeatedly.

\begin{successbox}
\textbf{Cost: \$0.} The IRS does not charge for an EIN. It is issued immediately online.\\[4pt]
\textbf{Warning:} Do NOT use third-party sites that charge \$50--200 to file your EIN. The IRS site is free.
\end{successbox}

\section{Write a Simple Operating Agreement}

An Operating Agreement is an internal document that describes how your LLC works. South Carolina does not require one, but having one:
\begin{itemize}[leftmargin=2em]
  \item Proves your LLC is a separate entity (protects your personal liability shield)
  \item Is required by banks and some payment processors
  \item Clarifies ownership if you ever add a partner
\end{itemize}

\step{9} Create a simple single-member Operating Agreement. You do not need a lawyer for this. A free template works:
\begin{actionbox}
Search for ``free single-member LLC operating agreement template'' on:
\begin{itemize}[leftmargin=1.5em]
  \item \url{https://www.northwestregisteredagent.com/free-llc-operating-agreement}
  \item \url{https://www.legalzoom.com/} (free template download)
\end{itemize}
Fill in your LLC name, your name as sole member, your address, and the date.\\
Sign it and keep a copy with your important business documents.
\end{actionbox}

\section{Check Local Business License Requirements}

\step{10} South Carolina does not require a statewide business license, but your city or county might. Check:
\begin{actionbox}
Contact your local county or city business license office. If you are in:\\[4pt]
\begin{tabular}{ll}
  Charleston County: & \url{https://www.charlestoncounty.org/departments/business-license/} \\
  Greenville County: & \url{https://www.greenvillecounty.org/AuditorBusinessLicense/} \\
  Columbia (Richland): & \url{https://www.richlandcountysc.gov/} \\
\end{tabular}\\[8pt]
Many SC counties charge \$0--50 for a basic home-based online business license. Some require no license at all for online-only businesses.
\end{actionbox}

\section{Chapter 1 Cost Summary}

\begin{table}[H]
\centering
\begin{tabular}{lr}
\toprule
\textbf{Item} & \textbf{Cost} \\
\midrule
Articles of Organization (online) & \$125.00 \\
EIN from IRS & \$0.00 \\
Operating Agreement (free template) & \$0.00 \\
Registered Agent (yourself) & \$0.00 \\
Local business license (estimated) & \$0.00--50.00 \\
\midrule
\textbf{Chapter 1 Total} & \textbf{\$125.00--175.00} \\
\bottomrule
\end{tabular}
\end{table}


% ══════════════════════════════════════════════════════════════
% CHAPTER 2: BANKING & FINANCES
% ══════════════════════════════════════════════════════════════
\chapter{Set Up Business Banking and Financial Accounts}

Never mix personal and business money. It destroys your liability protection and creates tax nightmares.

\section{Open a Business Checking Account}

\step{1} Gather your documents:
\begin{itemize}[leftmargin=2em]
  \item Articles of Organization (from Chapter 1)
  \item EIN confirmation letter (from Chapter 1)
  \item Your driver's license
  \item Operating Agreement
\end{itemize}

\step{2} Choose a bank. Recommended free options for small businesses:

\begin{table}[H]
\centering
\begin{tabularx}{\textwidth}{lXr}
\toprule
\textbf{Bank} & \textbf{Why} & \textbf{Monthly Fee} \\
\midrule
Mercury & Online-only, built for startups, easy Stripe integration & \$0 \\
Relay & Free business checking, up to 20 accounts & \$0 \\
Bluevine & 2\% interest on balances up to \$250K & \$0 \\
Local SC credit union & In-person support, community relationship & \$0--10 \\
\bottomrule
\end{tabularx}
\end{table}

\begin{infobox}
\textbf{Recommendation:} Open a \textbf{Mercury} account at \url{https://mercury.com}. It is 100\% online, free, integrates with Stripe, and is popular with solo founders. You can apply in 10 minutes.
\end{infobox}

\step{3} Deposit your startup capital. Transfer \$2{,}000 from your personal account into your new business account. This is your \textbf{capital contribution} to the LLC.

\section{Set Up Basic Bookkeeping}

\step{4} Sign up for \textbf{Wave Accounting} (free):
\begin{actionbox}
\url{https://www.waveapps.com}\\[4pt]
Wave is completely free accounting software. Connect your Mercury bank account so transactions import automatically. You will use this to:
\begin{itemize}[leftmargin=1.5em]
  \item Track all income and expenses
  \item Generate profit/loss reports
  \item Prepare for tax filing
\end{itemize}
\end{actionbox}

\step{5} Create expense categories in Wave:
\begin{multicols}{2}
\begin{itemize}[leftmargin=1.5em, itemsep=1pt]
  \item Hosting \& Infrastructure
  \item AI API Costs (OpenAI)
  \item Printing \& Fulfillment
  \item Shipping Costs
  \item Payment Processing Fees
  \item Domain \& Software
  \item Marketing \& Advertising
  \item Legal \& Professional
\end{itemize}
\end{multicols}

\section{Understand Your Tax Obligations}

\begin{warningbox}
\textbf{You must pay taxes.} A single-member LLC is a ``pass-through'' entity. All business profit is reported on your personal tax return.
\begin{itemize}[leftmargin=1.5em]
  \item \textbf{Federal:} Reported on Schedule C of your personal Form 1040
  \item \textbf{Self-employment tax:} 15.3\% on net profit (Social Security + Medicare)
  \item \textbf{South Carolina state income tax:} 0--6.5\% (graduated scale)
  \item \textbf{Estimated quarterly taxes:} If you expect to owe more than \$1{,}000, you must pay quarterly estimates (April 15, June 15, Sept 15, Jan 15)
  \item \textbf{Sales tax:} SC has a 6\% state sales tax plus county supplements (6--9\% total). Digital products may be exempt; physical books shipped in-state may be taxable. \textbf{Consult a CPA.}
\end{itemize}
\end{warningbox}

\step{6} Find a CPA. Budget approximately \$200--500/year for a CPA to file your taxes. Ask locally or search \url{https://www.thumbtack.com} for ``CPA small business South Carolina.'' This is money well spent.

\section{Chapter 2 Cost Summary}

\begin{table}[H]
\centering
\begin{tabular}{lr}
\toprule
\textbf{Item} & \textbf{Cost} \\
\midrule
Business checking account (Mercury) & \$0.00 \\
Bookkeeping software (Wave) & \$0.00 \\
CPA consultation (annual, estimated) & \$200--500 \\
\midrule
\textbf{Chapter 2 Total (upfront)} & \textbf{\$0.00} \\
\bottomrule
\end{tabular}
\end{table}

% ══════════════════════════════════════════════════════════════
% CHAPTER 3: DOMAIN & BRAND
% ══════════════════════════════════════════════════════════════
\chapter{Register Your Domain and Brand Assets}

\section{Register Your Domain Name}

\step{1} Go to Cloudflare Registrar:
\begin{actionbox}
\url{https://www.cloudflare.com/products/registrar/}\\[4pt]
Create a free Cloudflare account, then search for your desired \texttt{.com} domain.\\
Cloudflare charges \textbf{at-cost pricing} with zero markup: approximately \cost{\$10.44/year} for a \texttt{.com}. Free WHOIS privacy is included (hides your personal address from public records).
\end{actionbox}

\step{2} If your first choice is taken, try variations:
\begin{itemize}[leftmargin=2em]
  \item \texttt{crayonsandquills.com}
  \item \texttt{getcrayonsandquills.com}
  \item \texttt{crayons-and-quills.com}
  \item \texttt{mycrayonsandquills.com}
\end{itemize}

\section{Create a Simple Logo}

You do not need to hire a designer yet. Start with one of these free tools:

\step{3} Create a logo using one of:
\begin{table}[H]
\centering
\begin{tabularx}{\textwidth}{lXr}
\toprule
\textbf{Tool} & \textbf{Description} & \textbf{Cost} \\
\midrule
Canva & Drag-and-drop logo maker, thousands of templates & Free \\
Looka & AI-generated logo options & Free to browse, \$20 to download \\
Hatchful (Shopify) & Simple logo generator & Free \\
\bottomrule
\end{tabularx}
\end{table}

\begin{infobox}
\textbf{Recommendation:} Use \textbf{Canva} (\url{https://www.canva.com}). Search for ``children's book logo'' templates. Customize with your brand name and colors. Export as PNG (for web) and SVG (for print). You can always upgrade to a professional logo later.
\end{infobox}

\section{Establish Brand Colors and Fonts}

\step{4} For consistency, pick your brand colors now. Suggested palette:
\begin{itemize}[leftmargin=2em]
  \item \textbf{Primary:} A warm purple (\texttt{\#7E22CE}) --- trustworthy, creative
  \item \textbf{Accent:} Bright orange (\texttt{\#F97316}) --- playful, calls to action
  \item \textbf{Background:} Clean white with soft gray sections
  \item \textbf{Text:} Dark slate (\texttt{\#1E293B})
\end{itemize}

\noindent These colors are already configured in the Crayons \& Quills codebase.

\section{Chapter 3 Cost Summary}

\begin{table}[H]
\centering
\begin{tabular}{lr}
\toprule
\textbf{Item} & \textbf{Cost} \\
\midrule
Domain registration (Cloudflare, 1 year) & \$10.44 \\
Logo (Canva Free) & \$0.00 \\
\midrule
\textbf{Chapter 3 Total} & \textbf{\$10.44} \\
\bottomrule
\end{tabular}
\end{table}

% ══════════════════════════════════════════════════════════════
% CHAPTER 4: TECHNOLOGY ACCOUNTS
% ══════════════════════════════════════════════════════════════
\chapter{Set Up All Your Technology Accounts}

This chapter lists every account you need to create, in order. Do them all in one sitting---it takes about 2 hours. Keep a spreadsheet or password manager (like \textbf{Bitwarden}, free) with all your login credentials.

\begin{warningbox}
\textbf{Security:} Use a \textbf{unique, strong password} for each account. Install \textbf{Bitwarden} (free, \url{https://bitwarden.com}) to manage all passwords. Enable two-factor authentication (2FA) on every account that offers it.
\end{warningbox}

\section{Account Setup Checklist}

Work through this table from top to bottom. Each account takes 3--10 minutes.

\begin{longtable}{rp{1.8in}p{2.5in}r}
\toprule
\textbf{\#} & \textbf{Service} & \textbf{What It Does} & \textbf{Cost} \\
\midrule
\endhead
1 & Bitwarden & Password manager & Free \\
  & \multicolumn{3}{l}{\small\url{https://bitwarden.com} --- Store all your passwords here first} \\
\midrule
2 & OpenAI Platform & AI text and image generation & Pay-as-you-go \\
  & \multicolumn{3}{l}{\small\url{https://platform.openai.com} --- Add \$20 in credits to start} \\
\midrule
3 & Stripe & Payment processing & Free + fees \\
  & \multicolumn{3}{l}{\small\url{https://stripe.com} --- Requires EIN and business bank account} \\
\midrule
4 & Lulu Direct & Print-on-demand partner & Free \\
  & \multicolumn{3}{l}{\small\url{https://developers.lulu.com} --- Sign up for API access} \\
\midrule
5 & SendGrid & Transactional email & Free tier \\
  & \multicolumn{3}{l}{\small\url{https://sendgrid.com} --- Free: 100 emails/day} \\
\midrule
6 & Cloudflare & Domain \& DNS management & Free \\
  & \multicolumn{3}{l}{\small\url{https://cloudflare.com} --- Already created for domain registration} \\
\midrule
7 & Railway.app & Application hosting & \$5--20/mo \\
  & \multicolumn{3}{l}{\small\url{https://railway.app} --- Where your website runs} \\
\midrule
8 & GitHub & Code storage & Free \\
  & \multicolumn{3}{l}{\small\url{https://github.com} --- Where your code lives} \\
\midrule
9 & Claude Code (Claude Max) & AI coding assistant & \$0 (included) \\
  & \multicolumn{3}{l}{\small\url{https://claude.ai/code} --- You already have Claude Max. Use Claude Code to build/modify the app.} \\
\midrule
10 & Make.com & Automation workflows & \$10.59/mo \\
   & \multicolumn{3}{l}{\small\url{https://make.com} --- Connects everything together automatically} \\
\midrule
11 & Google Workspace or Gmail & Business email & Free \\
   & \multicolumn{3}{l}{\small Use a dedicated Gmail like \texttt{crayonsandquills@gmail.com} to start} \\
\bottomrule
\end{longtable}

\section{Detailed Setup: OpenAI Platform}

This is the ``brain'' of your business---it generates stories and illustrations.

\step{5} Go to \url{https://platform.openai.com} and create an account.

\step{6} Navigate to \textbf{Settings} $\rightarrow$ \textbf{Billing}. Add a payment method (your business debit card).

\step{7} Add \cost{\$20} in prepaid credits. This is enough to generate approximately 20--25 complete books during testing.

\step{8} Navigate to \textbf{API Keys}. Click \textbf{Create new secret key}. Name it ``Crayons \& Quills Production.'' \textbf{Copy this key immediately}---you will never see it again. Save it in Bitwarden.

\begin{warningbox}
\textbf{Never share your API key.} Never paste it into a public website, email, or chat. Anyone with this key can run up charges on your account. Set a monthly spending limit under \textbf{Settings} $\rightarrow$ \textbf{Limits} to protect yourself. Set it to \cost{\$50/month} initially.
\end{warningbox}

\section{Detailed Setup: Stripe}

\step{9} Go to \url{https://stripe.com} and click \textbf{Start now}.

\step{10} Complete business verification:
\begin{itemize}[leftmargin=2em]
  \item Business type: \textbf{LLC}
  \item Business name: Your LLC name
  \item EIN: Enter your EIN from Chapter 1
  \item Business address: Your SC address
  \item Bank account: Connect your Mercury business account
  \item Identity verification: Upload your driver's license
\end{itemize}

\step{11} Once approved (usually same day), go to \textbf{Developers} $\rightarrow$ \textbf{API Keys}. You will see:
\begin{itemize}[leftmargin=2em]
  \item \textbf{Publishable key} (starts with \texttt{pk\_live\_}) --- safe to use in frontend code
  \item \textbf{Secret key} (starts with \texttt{sk\_live\_}) --- \textbf{keep secret}, save in Bitwarden
\end{itemize}

\step{12} Create your products in Stripe Dashboard $\rightarrow$ \textbf{Products}:

\begin{infobox}
\textbf{Dynamic pricing:} Your backend calculates the exact price based on each customer's selections (page count, binding, paper, size, and add-ons). The Stripe payment intent is created dynamically---you do \textbf{not} need a fixed Stripe Product for every combination. However, create these base products for your dashboard and subscription tracking:
\end{infobox}

\begin{table}[H]
\centering
\begin{tabularx}{\textwidth}{Xlr}
\toprule
\textbf{Product Name} & \textbf{Type} & \textbf{Starting Price} \\
\midrule
Personalized Children's Book (12-page) & One-time & From \$19.99 \\
Personalized Children's Book (24-page) & One-time & From \$24.99 \\
Personalized Children's Book (36-page) & One-time & From \$34.99 \\
Hardcover Upgrade & Add-on & +\$10.00 \\
Premium Matte Paper & Add-on & +\$5.00 \\
Glossy Paper \& Cover & Add-on & +\$8.00 \\
Larger Size (8.5$\times$11 or 11$\times$8.5) & Add-on & +\$3.00 \\
Birthday Club Subscription & Recurring (yearly) & \$24.99/yr \\
Adventure Series Subscription & Recurring (quarterly) & \$22.99/qtr \\
Digital PDF Add-on & One-time & +\$4.99 \\
Audiobook Add-on & One-time & +\$9.99 \\
Gift Wrap & One-time & +\$3.99 \\
\bottomrule
\end{tabularx}
\end{table}

\section{Detailed Setup: Lulu Direct}

\step{13} Go to \url{https://developers.lulu.com} and create a developer account.

\step{14} From the developer dashboard, create an \textbf{API Application}. You will receive:
\begin{itemize}[leftmargin=2em]
  \item \textbf{Client Key} (like a username for the API)
  \item \textbf{Client Secret} (like a password) --- save in Bitwarden
\end{itemize}

\step{15} Familiarize yourself with Lulu's pricing calculator:
\begin{actionbox}
\url{https://www.lulu.com/pricing}\\[4pt]
Enter these specs to see your per-book cost:
\begin{itemize}[leftmargin=1.5em]
  \item Trim size: 8.5\,'' $\times$ 8.5\,''
  \item Page count: 24 (or 50)
  \item Interior: Full color
  \item Paper: Standard 60\#
  \item Binding: Saddle-stitch (under 32 pages) or Perfect-bound (32+ pages)
  \item Cover: Softcover
\end{itemize}
\end{actionbox}

\begin{infobox}
\textbf{Important:} Books under 32 pages are \textbf{saddle-stitched} (stapled) on Lulu, not perfect-bound. If you want a spine (perfect binding), your minimum is 32 pages. A 24-page saddle-stitched book still looks professional---many children's books use this format.
\end{infobox}

\section{Detailed Setup: Railway.app}

\step{16} Go to \url{https://railway.app} and sign up with your GitHub account.

\step{17} Choose the \textbf{Hobby plan} at \cost{\$5/month} (includes \$5 of usage credits, which covers most small projects). Upgrade to Pro (\$20/month) when you start getting consistent orders.

\step{18} You will deploy three things on Railway later (Chapter 6):
\begin{enumerate}[leftmargin=2em]
  \item Your backend API (NestJS)
  \item Your frontend website (Next.js)
  \item A PostgreSQL database
  \item A Redis instance (for background job queues)
\end{enumerate}

\section{Chapter 4 Cost Summary}

\begin{table}[H]
\centering
\begin{tabular}{lr}
\toprule
\textbf{Item} & \textbf{Cost} \\
\midrule
Bitwarden (password manager) & \$0.00 \\
OpenAI initial credits & \$20.00 \\
Stripe account & \$0.00 \\
Lulu Direct account & \$0.00 \\
SendGrid account (free tier) & \$0.00 \\
Cloudflare account & \$0.00 \\
Railway.app (first month) & \$5.00 \\
GitHub account & \$0.00 \\
Claude Code (via Claude Max) & \$0.00 (already subscribed) \\
Make.com Core (first month) & \$10.59 \\
\midrule
\textbf{Chapter 4 Total} & \textbf{\$35.59} \\
\bottomrule
\end{tabular}
\end{table}

% ══════════════════════════════════════════════════════════════
% CHAPTER 5: CONFIGURE AND RUN THE APP
% ══════════════════════════════════════════════════════════════
\chapter{Configure and Run the Application Using Claude Code}

Your application is \textbf{already built}. The \texttt{LLM-books} repository contains a complete Next.js frontend, NestJS backend, Prisma database schema, AI service layer, background job workers, and all the integration code for Stripe, SendGrid, and Lulu Direct. You do not need to write code from scratch.

What you need to do now is \textbf{configure} it (add your API keys), \textbf{test} it (run it locally or in the cloud), and \textbf{deploy} it (make it live on the internet). You will use \textbf{Claude Code} for all of this---through plain English conversation.

\begin{successbox}
\textbf{Advantage of Claude Max:} Because you already pay for Claude Max, your AI coding tool costs \textbf{\$0 extra per month}. Other founders pay \$20--40/month for Cursor, Windsurf, or Bolt.new. You save that entirely.
\end{successbox}

\section{What Already Exists in Your Codebase}

Before you begin, understand what you already have:

\begin{table}[H]
\centering
\small
\begin{tabularx}{\textwidth}{lX}
\toprule
\textbf{Directory / File} & \textbf{What It Contains} \\
\midrule
\texttt{backend/src/auth/} & JWT authentication with refresh tokens, login/register endpoints \\
\texttt{backend/src/books/} & Book CRUD endpoints and generation trigger \\
\texttt{backend/src/ai/story/} & StoryEngineService --- generates age-appropriate stories via GPT-4o \\
\texttt{backend/src/ai/illustration/} & IllustrationEngineService --- generates images via DALL-E 3 \\
\texttt{backend/src/ai/character/} & CharacterEngineService --- GPT-4 Vision character extraction \\
\texttt{backend/src/ai/pdf/} & PdfAssemblyService --- print-ready and digital PDF generation \\
\texttt{backend/src/payments/} & Stripe integration with webhook handling \\
\texttt{backend/src/fulfillment/} & Lulu Direct API integration \\
\texttt{backend/src/email/} & SendGrid email templates (order confirmation, shipping, etc.) \\
\texttt{backend/prisma/schema.prisma} & Database schema with 14 models, indexes, and enums \\
\texttt{frontend/src/app/book/create/} & Multi-step book creation wizard \\
\texttt{frontend/src/app/dashboard/} & User dashboard for viewing orders and books \\
\texttt{frontend/src/app/(marketing)/} & SEO-optimized landing pages \\
\texttt{agents/} & Six Mac Mini agent script stubs (order monitor, fulfillment, etc.) \\
\texttt{CLAUDE.md} & Persistent instructions for Claude Code sessions \\
\texttt{.claude/settings.json} & SessionStart hook for automatic project status display \\
\bottomrule
\end{tabularx}
\end{table}

\section{The Process From Here}

Your remaining steps are:

\begin{enumerate}[leftmargin=2em]
  \item Open the codebase in \textbf{Claude Code} (web or CLI)
  \item Let Claude install dependencies (\texttt{npm install} in both \texttt{backend/} and \texttt{frontend/})
  \item Configure your \texttt{.env} files with the API keys from Chapter 4
  \item Run the database migration (\texttt{npx prisma migrate dev})
  \item Test locally (or deploy directly to Railway---see Chapter 6)
  \item Use Claude Code for all future edits, fixes, and feature additions
\end{enumerate}

\section{Using Claude Code on the Web (Easiest)}

You can use Claude Code directly in your browser---no software installation required on your computer.

\step{1} Push your code to GitHub (if not already there):
\begin{actionbox}
Go to \url{https://github.com} and create a new repository called \texttt{LLM-books}.\\
Upload the codebase, or use Claude Code to help you push it.
\end{actionbox}

\step{2} Open Claude Code in your browser:
\begin{actionbox}
Go to \url{https://claude.ai/code}\\[4pt]
Your Claude Max subscription gives you full access. Claude Code on the web runs in a cloud environment with a terminal, file editor, and Git---everything you need.
\end{actionbox}

\step{3} Clone and set up the project:
\begin{actionbox}
In the Claude Code chat, type:\\[4pt]
\textit{``Clone my repository at https://github.com/YourUsername/LLM-books.git, install all dependencies for both the backend and frontend, and copy the .env.example files to create my local config files.''}\\[4pt]
Claude will run all the terminal commands for you automatically. Specifically, it will run:
\begin{itemize}[leftmargin=1.5em, itemsep=2pt]
  \item \texttt{git clone} your repository
  \item \texttt{npm install} in both \texttt{backend/} and \texttt{frontend/}
  \item \texttt{cp backend/.env.example backend/.env}
  \item \texttt{cp frontend/.env.example frontend/.env.local}
\end{itemize}
\end{actionbox}

\begin{infobox}
\textbf{What happens on session start:} Because your repository includes a \texttt{.claude/settings.json} file with a SessionStart hook, Claude Code will automatically print a project overview showing the tech stack, key directories, current git branch, and dependency status every time you open a session. This is already configured---you do not need to set it up.
\end{infobox}

\step{4} Configure your API keys:
\begin{actionbox}
Your \texttt{backend/.env.example} already lists every variable you need. Tell Claude:\\[4pt]
\textit{``Help me fill in my backend/.env file. Here are my keys:}\\
\textit{OpenAI: sk-...}\\
\textit{Stripe Secret: sk\_live\_...}\\
\textit{Stripe Publishable: pk\_live\_...}\\
\textit{SendGrid: SG...}\\
\textit{Lulu API Key: ...}\\
\textit{Lulu API Secret: ...''}\\[4pt]
Claude will edit the \texttt{.env} file with your actual keys. The file already contains placeholder values for:\\
\texttt{DATABASE\_URL}, \texttt{REDIS\_URL}, \texttt{JWT\_SECRET}, \texttt{OPENAI\_API\_KEY}, \texttt{STRIPE\_SECRET\_KEY}, \texttt{STRIPE\_WEBHOOK\_SECRET}, \texttt{SENDGRID\_API\_KEY}, \texttt{LULU\_API\_KEY}, \texttt{LULU\_API\_SECRET}, \texttt{S3\_BUCKET\_NAME}, and more.
\end{actionbox}

\begin{warningbox}
\textbf{Security note:} Claude Code conversations are private to your account, but as a best practice, avoid pasting API keys into any tool you do not fully trust. Claude Code runs in your own sandboxed environment and does not share your data.
\end{warningbox}

\step{5} Set up the database and test locally:
\begin{actionbox}
Tell Claude:\\[4pt]
\textit{``Run the Prisma migration to create the database tables, then start the backend and frontend dev servers so I can test locally.''}\\[4pt]
Claude will run:
\begin{itemize}[leftmargin=1.5em, itemsep=2pt]
  \item \texttt{npx prisma generate} (generates the Prisma client from your schema)
  \item \texttt{npx prisma migrate dev --name init} (creates all 14 database tables)
  \item \texttt{npm run start:dev} in \texttt{backend/} (starts the API on port 4000)
  \item \texttt{npm run dev} in \texttt{frontend/} (starts the website on port 3000)
\end{itemize}
You can then visit \texttt{http://localhost:3000} to see your Crayons \& Quills site running.
\end{actionbox}

\step{6} Set up the frontend config:
\begin{actionbox}
\textit{``Now set up the frontend/.env.local file with my Stripe publishable key pk\_live\_... and set the API URL to my backend.''}\\[4pt]
The \texttt{frontend/.env.example} already lists the required variables: \texttt{NEXT\_PUBLIC\_API\_URL} and \texttt{NEXT\_PUBLIC\_STRIPE\_PUBLISHABLE\_KEY}.
\end{actionbox}

\section{Using Claude Code Locally (Alternative)}

If you prefer to run everything on your own computer:

\step{1} Install \textbf{Git} (the tool that manages code versions):
\begin{actionbox}
\textbf{On Windows:} Download from \url{https://git-scm.com/download/win} and run the installer. Accept all defaults.\\[4pt]
\textbf{On Mac:} Open Terminal (search for it in Spotlight) and type \texttt{git --version}. If not installed, it will prompt you to install it.
\end{actionbox}

\step{2} Install \textbf{Node.js} (the engine that runs the application):
\begin{actionbox}
\url{https://nodejs.org}\\[4pt]
Download the \textbf{LTS} version (not ``Current''). Run the installer. Accept all defaults.
\end{actionbox}

\step{3} Install \textbf{Claude Code CLI}:
\begin{actionbox}
Open a terminal and run:\\[4pt]
\texttt{npm install -g @anthropic-ai/claude-code}\\[4pt]
Then navigate to your project folder and run:\\[4pt]
\texttt{cd LLM-books}\\
\texttt{claude}\\[4pt]
This opens Claude Code in your terminal, connected to your codebase.
\end{actionbox}

\step{4} Let Claude handle setup:
\begin{actionbox}
In the Claude Code prompt, type:\\[4pt]
\textit{``Install all dependencies for both backend and frontend, copy .env.example files to .env, run the Prisma migration, and help me fill in my API keys.''}\\[4pt]
Claude runs \texttt{npm install} in both directories, copies the environment files, runs \texttt{npx prisma migrate dev}, and walks you through entering your API keys from Chapter 4. Your existing \texttt{CLAUDE.md} file gives Claude the context it needs to understand the project structure automatically.
\end{actionbox}

\section{How to Make Changes Without Coding}

Your codebase is fully structured with clear module boundaries. Whenever you need to change something, open Claude Code and describe what you want. Claude reads your \texttt{CLAUDE.md} file automatically and understands the entire project:

\begin{infobox}
\textbf{Example prompts you can type into Claude Code:}
\begin{itemize}[leftmargin=1.5em, itemsep=4pt]
  \item ``Change the price of the 24-page book from \$29.99 to \$34.99'' \\ {\small(Claude edits the Stripe product config and any hardcoded references)}
  \item ``Add a new landing page for books about unicorns'' \\ {\small(Claude creates a new page in \texttt{frontend/src/app/(marketing)/} following the existing pattern)}
  \item ``The checkout page is showing an error---help me fix it'' \\ {\small(Claude reads error logs, traces the issue through the frontend/backend, and patches it)}
  \item ``Add a field for the child's school name in the wizard'' \\ {\small(Claude updates the Prisma schema, generates a migration, and updates the wizard component)}
  \item ``Run the tests and fix any failures''
  \item ``Deploy the latest changes to Railway''
\end{itemize}
Claude reads your entire codebase, makes the changes, and can commit and push them to GitHub for you.
\end{infobox}

\begin{successbox}
\textbf{Already configured:} Your repository includes a \textbf{SessionStart hook} in \texttt{.claude/settings.json} that prints a project status report every time you start a Claude Code session. It also includes \texttt{CLAUDE.md} with persistent project instructions, so Claude always knows your tech stack, key directories, database setup, and operational rules.
\end{successbox}

\section{Chapter 5 Cost Summary}

\begin{table}[H]
\centering
\begin{tabular}{lr}
\toprule
\textbf{Item} & \textbf{Cost} \\
\midrule
Claude Code (via Claude Max) & \$0.00 (already subscribed) \\
Node.js & \$0.00 \\
Git & \$0.00 \\
\midrule
\textbf{Chapter 5 Total} & \textbf{\$0.00} \\
\bottomrule
\end{tabular}
\end{table}


% ══════════════════════════════════════════════════════════════
% CHAPTER 6: DEPLOY
% ══════════════════════════════════════════════════════════════
\chapter{Deploy the Application (Make It Live on the Internet)}

``Deploying'' means putting your application on a server so anyone in the world can visit your website. Your codebase is already structured for deployment---it includes separate \texttt{backend/} and \texttt{frontend/} directories with their own \texttt{package.json} files, a \texttt{docker-compose.yml} for production, and a Prisma migration system. We use Railway.app because it is the simplest option for non-technical founders.

\section{What Gets Deployed}

Your repository already contains all four components. Each deploys separately on Railway:

\begin{table}[H]
\centering
\begin{tabularx}{\textwidth}{lX}
\toprule
\textbf{Component} & \textbf{What It Does} \\
\midrule
Frontend (Next.js) & The website customers see and interact with \\
Backend (NestJS) & The ``brain'' that processes orders, talks to OpenAI, etc. \\
PostgreSQL Database & Stores all data: users, books, orders, payments \\
Redis & Manages the queue of book generation jobs \\
\bottomrule
\end{tabularx}
\end{table}

\section{Deploy to Railway.app}

\step{1} Log in to \url{https://railway.app} with your GitHub account.

\step{2} Create a new project. Click \textbf{New Project} $\rightarrow$ \textbf{Deploy from GitHub Repo} $\rightarrow$ Select your \texttt{LLM-books} repository. Railway will detect your project structure automatically.

\step{3} Add a PostgreSQL database:
\begin{actionbox}
In your Railway project, click \textbf{New} $\rightarrow$ \textbf{Database} $\rightarrow$ \textbf{PostgreSQL}.\\
Railway provisions it instantly. Click on the database and go to \textbf{Variables} tab. Copy the \texttt{DATABASE\_URL} value---you will need it.
\end{actionbox}

\step{4} Add Redis:
\begin{actionbox}
Click \textbf{New} $\rightarrow$ \textbf{Database} $\rightarrow$ \textbf{Redis}.\\
Copy the \texttt{REDIS\_URL} from the Variables tab.
\end{actionbox}

\step{5} Configure the backend service:
\begin{actionbox}
Click on your deployed service. Go to \textbf{Settings}:
\begin{itemize}[leftmargin=1.5em]
  \item Root Directory: \texttt{backend}
  \item Build Command: \texttt{npm install \&\& npx prisma generate \&\& npm run build}
  \item Start Command: \texttt{npx prisma migrate deploy \&\& npm run start:prod}
\end{itemize}
These commands work because your \texttt{backend/package.json} already defines the \texttt{build} and \texttt{start:prod} scripts, and \texttt{backend/prisma/schema.prisma} contains the full database schema.\\[4pt]
Go to \textbf{Variables} and add ALL the environment variables from your \texttt{backend/.env} file. Railway has a bulk import feature---paste the entire \texttt{.env} file contents. Your \texttt{backend/.env.example} lists every required variable.
\end{actionbox}

\step{6} Deploy the frontend as a separate service:
\begin{actionbox}
Click \textbf{New} $\rightarrow$ \textbf{GitHub Repo} $\rightarrow$ Select the same repository.\\
In Settings:
\begin{itemize}[leftmargin=1.5em]
  \item Root Directory: \texttt{frontend}
  \item Build Command: \texttt{npm install \&\& npm run build}
  \item Start Command: \texttt{npm start}
\end{itemize}
Add the frontend environment variables (\texttt{NEXT\_PUBLIC\_API\_URL}, etc.).
\end{actionbox}

\step{7} Set up custom domains:
\begin{actionbox}
In Railway, click on each service $\rightarrow$ \textbf{Settings} $\rightarrow$ \textbf{Networking} $\rightarrow$ \textbf{Custom Domain}.\\[4pt]
\begin{itemize}[leftmargin=1.5em]
  \item Frontend: \texttt{www.yourdomain.com} and \texttt{yourdomain.com}
  \item Backend: \texttt{api.yourdomain.com}
\end{itemize}
Railway gives you DNS records to add in Cloudflare. Go to your Cloudflare dashboard $\rightarrow$ DNS $\rightarrow$ Add the CNAME records Railway provides.
\end{actionbox}

\step{8} Enable HTTPS. Railway provides free SSL/TLS certificates automatically. No action needed.

\section{Verify Deployment}

\step{9} Visit your domain in a browser. You should see your Crayons \& Quills homepage. Test:
\begin{itemize}[leftmargin=2em]
  \item[$\square$] Homepage loads with hero section
  \item[$\square$] Navigation links work
  \item[$\square$] Book creation wizard loads
  \item[$\square$] User registration works
  \item[$\square$] Login works
\end{itemize}

\section{Estimated Monthly Hosting Costs}

\begin{table}[H]
\centering
\begin{tabular}{lr}
\toprule
\textbf{Service} & \textbf{Estimated Monthly Cost} \\
\midrule
Railway Hobby plan base & \$5.00 \\
Backend compute (light usage) & \$3--8.00 \\
Frontend compute & \$2--5.00 \\
PostgreSQL (small) & \$1--5.00 \\
Redis (small) & \$1--3.00 \\
\midrule
\textbf{Total estimated hosting} & \textbf{\$12--26/month} \\
\bottomrule
\end{tabular}
\end{table}

\begin{infobox}
Railway uses \textbf{per-minute billing}---you only pay for what you use. A low-traffic site in the early days may cost as little as \$8--12/month total.
\end{infobox}

% ══════════════════════════════════════════════════════════════
% CHAPTER 7: AI AGENT AUTOMATION
% ══════════════════════════════════════════════════════════════
\chapter{Set Up AI Agent Automation (The ``Hands-Off'' Layer)}

This is the chapter that makes your business run itself. Your backend already has all the API endpoints for book generation, fulfillment, and email notifications (see \texttt{docs/API\_ROUTES.md} in your repository). Now you need to \textbf{connect} them so they trigger automatically.

You will use \textbf{Make.com} to create workflows that call your existing API. Later (Chapter 8), you can migrate these to your Mac Mini for more power and lower cost.

\section{What Is Make.com?}

Make.com is a visual automation tool. You connect ``modules'' (like Lego blocks) to create workflows that call your backend's existing REST API:

\begin{center}
\textit{When a new order is paid (Stripe)} $\rightarrow$ \textit{Call your API to generate the book} $\rightarrow$ \textit{Call your API to send to printer} $\rightarrow$ \textit{Email the customer}
\end{center}

Your backend already handles all the complex work (AI generation via \texttt{backend/src/ai/}, PDF assembly, Lulu submission via \texttt{backend/src/fulfillment/}). Make.com simply triggers the right endpoints at the right time. You set this up once, and it runs automatically.

\section{Workflow 1: New Order Processing}

This is your most important automation. When a customer pays, everything happens automatically.

\step{1} Log in to \url{https://make.com}. Click \textbf{Create a new scenario}.

\step{2} Build this workflow:

\begin{table}[H]
\centering
\small
\begin{tabularx}{\textwidth}{rlX}
\toprule
\textbf{Step} & \textbf{Module} & \textbf{What It Does} \\
\midrule
1 & Stripe: Watch Events & Triggers when \texttt{payment\_intent.succeeded} fires \\
2 & HTTP: POST & Calls your backend API: \texttt{POST /api/v1/books/\{id\}/generate} \\
3 & Wait / Delay & Waits 5--15 minutes for AI generation to complete \\
4 & HTTP: GET & Checks book status: \texttt{GET /api/v1/books/\{id\}} \\
5 & Router & If status = REVIEW\_READY $\rightarrow$ continue. If FAILED $\rightarrow$ alert you \\
6 & HTTP: POST & Submits to Lulu: \texttt{POST /api/v1/fulfillment/orders/\{id\}/submit} \\
7 & SendGrid: Send Email & Sends order confirmation to customer \\
\bottomrule
\end{tabularx}
\end{table}

\begin{infobox}
\textbf{For non-technical users:} Make.com has a visual drag-and-drop interface. You do not type code. You click on a module (e.g., ``Stripe''), select the action (``Watch Events''), and fill in your API key. Make.com handles the rest. Their documentation includes video tutorials for every module.
\end{infobox}

\section{Workflow 2: Fulfillment Status Tracking}

This workflow checks Lulu every few hours for shipping updates and notifies customers.

\begin{table}[H]
\centering
\small
\begin{tabularx}{\textwidth}{rlX}
\toprule
\textbf{Step} & \textbf{Module} & \textbf{What It Does} \\
\midrule
1 & Schedule & Runs every 4 hours \\
2 & HTTP: GET & Calls your backend: \texttt{GET /api/v1/fulfillment/orders/\{id\}/status} \\
3 & Filter & Only continue if status changed (e.g., now ``shipped'') \\
4 & SendGrid: Send Email & Sends shipping notification with tracking link \\
\bottomrule
\end{tabularx}
\end{table}

\section{Workflow 3: Subscription Renewal Processing}

\begin{table}[H]
\centering
\small
\begin{tabularx}{\textwidth}{rlX}
\toprule
\textbf{Step} & \textbf{Module} & \textbf{What It Does} \\
\midrule
1 & Stripe: Watch Events & Triggers on \texttt{invoice.payment\_succeeded} for subscriptions \\
2 & HTTP: POST & Creates a new book for the child (updated age) \\
3 & HTTP: POST & Starts generation pipeline \\
4 & SendGrid: Send Email & ``Your child's new book is being created!'' \\
\bottomrule
\end{tabularx}
\end{table}

\section{Workflow 4: Failed Generation Alert}

\begin{table}[H]
\centering
\small
\begin{tabularx}{\textwidth}{rlX}
\toprule
\textbf{Step} & \textbf{Module} & \textbf{What It Does} \\
\midrule
1 & Schedule & Runs every 30 minutes \\
2 & HTTP: GET & Checks for books with status = FAILED \\
3 & Filter & Only continue if there are failed books \\
4 & Gmail: Send Email & Sends YOU an alert: ``A book generation failed---review needed'' \\
\bottomrule
\end{tabularx}
\end{table}

\section{Workflow 5: Abandoned Cart Recovery}

\begin{table}[H]
\centering
\small
\begin{tabularx}{\textwidth}{rlX}
\toprule
\textbf{Step} & \textbf{Module} & \textbf{What It Does} \\
\midrule
1 & Schedule & Runs daily at 10 AM \\
2 & HTTP: GET & Finds books in DRAFT status older than 24 hours \\
3 & Filter & Only users who haven't been emailed already \\
4 & SendGrid: Send Email & ``You started creating a book---finish it!'' with a link back \\
\bottomrule
\end{tabularx}
\end{table}

\section{Make.com Costs at Scale}

\begin{table}[H]
\centering
\begin{tabular}{lrr}
\toprule
\textbf{Monthly Orders} & \textbf{Est.\ Operations} & \textbf{Make.com Plan Needed} \\
\midrule
1--20 & $\sim$500 & Core (\$10.59/mo) \\
20--100 & $\sim$3{,}000 & Core (\$10.59/mo) \\
100--500 & $\sim$15{,}000 & Pro (\$18.82/mo) \\
500+ & $\sim$50{,}000+ & Teams (\$29/mo) \\
\bottomrule
\end{tabular}
\end{table}


% ══════════════════════════════════════════════════════════════
% CHAPTER 8: MAC MINI AGENT RUNNER
% ══════════════════════════════════════════════════════════════
\chapter{Set Up Your Mac Mini M2 as an Always-On Agent Runner}

You own a \textbf{Mac Mini M2}---this is a significant advantage. Instead of relying entirely on Make.com (\$11--29/month) for automation, you can run AI agents \textbf{directly on your Mac Mini} as an always-on business operations machine. This gives you more power, more flexibility, and lower costs than any cloud automation platform.

\section{Why a Dedicated Agent Machine?}

\begin{table}[H]
\centering
\begin{tabularx}{\textwidth}{lXX}
\toprule
& \textbf{Make.com (Cloud)} & \textbf{Mac Mini M2 (Local)} \\
\midrule
Monthly cost & \$11--29/mo & \$0 (electricity only) \\
Complexity limit & Simple linear workflows & Arbitrarily complex AI agents \\
AI reasoning & None (just HTTP calls) & Full Claude/GPT reasoning \\
Custom logic & Limited & Unlimited (any code) \\
Uptime & 99.9\% & 99\%+ (with auto-restart) \\
Data privacy & Third-party servers & Your own hardware \\
\bottomrule
\end{tabularx}
\end{table}

\begin{successbox}
\textbf{Recommended hybrid approach:} Use Make.com for the simplest webhook-triggered workflows (payment $\rightarrow$ trigger generation) and your Mac Mini for complex agent tasks (monitoring, recovery, customer communication, analytics). As you gain confidence, migrate everything to the Mac Mini and cancel Make.com entirely.
\end{successbox}

\section{Initial Setup}

\step{1} Configure your Mac Mini as an always-on server:
\begin{actionbox}
\textbf{System Settings} $\rightarrow$ \textbf{Energy}:\\
\begin{itemize}[leftmargin=1.5em]
  \item Enable ``Start up automatically after a power failure''
  \item Set ``Turn display off after'' to 1 minute (saves energy, Mac keeps running)
  \item Disable sleep: ``Prevent automatic sleeping when the display is off'' $\rightarrow$ ON
\end{itemize}
\textbf{System Settings} $\rightarrow$ \textbf{General} $\rightarrow$ \textbf{Login Items \& Extensions}:\\
\begin{itemize}[leftmargin=1.5em]
  \item Add your agent scripts to ``Open at Login'' (we will create these next)
\end{itemize}
\end{actionbox}

\step{2} Install the core tools:
\begin{actionbox}
Open Terminal on the Mac Mini and run:\\[4pt]
\texttt{/bin/bash -c "\$(curl -fsSL https://raw.githubusercontent.com/Homebrew/install/HEAD/install.sh)"}\\[4pt]
Then install everything you need:\\[4pt]
\texttt{brew install node git}\\
\texttt{npm install -g @anthropic-ai/claude-code}\\
\texttt{npm install -g pm2}\\[4pt]
\textbf{pm2} is a process manager that keeps your agent scripts running 24/7, restarts them if they crash, and provides logs.
\end{actionbox}

\step{3} Clone your codebase on the Mac Mini:
\begin{actionbox}
\texttt{cd \~{}}\\
\texttt{git clone https://github.com/YourUsername/LLM-books.git}\\
\texttt{cd LLM-books}\\[4pt]
Your repository already includes the \texttt{agents/} directory with stub scripts, \texttt{CLAUDE.md} for Claude Code context, and \texttt{.claude/settings.json} for automatic session startup.
\end{actionbox}

\section{Agent Architecture}

Your Mac Mini will run several lightweight agent scripts, each handling a specific business function:

\begin{table}[H]
\centering
\small
\begin{tabularx}{\textwidth}{lXl}
\toprule
\textbf{Agent} & \textbf{Responsibility} & \textbf{Runs} \\
\midrule
Order Monitor & Watches for new paid orders, triggers book generation & Continuous \\
Fulfillment Tracker & Polls Lulu for shipping status updates & Every 4 hours \\
Failed Job Recovery & Detects and retries failed book generations & Every 30 min \\
Abandoned Cart Agent & Sends reminder emails for unfinished books & Daily 10 AM \\
Health Monitor & Checks that all services (Railway, APIs) are up & Every 15 min \\
Daily Digest Agent & Emails you a summary of orders, revenue, issues & Daily 8 AM \\
\bottomrule
\end{tabularx}
\end{table}

\section{Completing the Agent Scripts}

Your repository already contains an \texttt{agents/} directory with \textbf{stub files} for all six agents:

\begin{itemize}[leftmargin=2em, itemsep=2pt]
  \item \texttt{agents/order-monitor.ts}
  \item \texttt{agents/fulfillment-tracker.ts}
  \item \texttt{agents/failed-job-recovery.ts}
  \item \texttt{agents/abandoned-cart.ts}
  \item \texttt{agents/health-monitor.ts}
  \item \texttt{agents/daily-digest.ts}
\end{itemize}

These stubs define the structure and purpose of each agent. You need to \textbf{flesh them out} with your specific API URL and configuration. Use Claude Code to do this automatically.

\step{4} Open Claude Code on the Mac Mini:
\begin{actionbox}
\texttt{cd \~{}/LLM-books}\\
\texttt{claude}\\[4pt]
Then tell Claude:\\[4pt]
\textit{``I already have agent stubs in the agents/ directory. Flesh out each one with full implementation:
\begin{enumerate}[leftmargin=1.5em]
  \item order-monitor.ts --- polls my backend API at [your Railway URL] every 60 seconds for new paid orders and triggers book generation via POST /api/v1/books/\{id\}/generate
  \item fulfillment-tracker.ts --- every 4 hours, checks all orders with status SENT\_TO\_PRINT and updates their fulfillment status from Lulu
  \item failed-job-recovery.ts --- every 30 minutes, finds books with status FAILED and retries generation up to 3 times
  \item abandoned-cart.ts --- daily at 10 AM, finds books in DRAFT status older than 24 hours and sends reminder emails
  \item health-monitor.ts --- every 15 minutes, pings the backend API, OpenAI, and Stripe to verify they respond
  \item daily-digest.ts --- every day at 8 AM, compiles yesterday's order count, revenue, and issues, and emails me a summary
\end{enumerate}
Each script should use environment variables from .env, include proper error handling and logging, and be designed to run under pm2.''}
\end{actionbox}

\step{5} Claude Code will complete all six scripts based on the existing stubs. Review them, then start them with pm2:
\begin{actionbox}
\texttt{pm2 start agents/order-monitor.ts --name order-monitor}\\
\texttt{pm2 start agents/fulfillment-tracker.ts --name fulfillment-tracker}\\
\texttt{pm2 start agents/failed-job-recovery.ts --name failed-recovery}\\
\texttt{pm2 start agents/abandoned-cart.ts --name abandoned-cart}\\
\texttt{pm2 start agents/health-monitor.ts --name health-monitor}\\
\texttt{pm2 start agents/daily-digest.ts --name daily-digest}\\[4pt]
Save the process list so they survive reboots:\\[4pt]
\texttt{pm2 save}\\
\texttt{pm2 startup}\\[4pt]
pm2 will give you a command to copy/paste---run it to enable auto-start on boot.
\end{actionbox}

\section{Monitoring Your Agents}

\step{6} Check agent status any time:
\begin{actionbox}
\texttt{pm2 status} \quad--- shows all running agents, uptime, restarts, CPU/memory\\
\texttt{pm2 logs} \quad--- live stream of all agent logs\\
\texttt{pm2 logs order-monitor} \quad--- logs for a specific agent\\
\texttt{pm2 monit} \quad--- real-time dashboard
\end{actionbox}

\section{Advanced: Claude Agent SDK (Optional)}

For even more powerful agents that can \textit{reason} about problems (not just follow rules), you can use the \textbf{Claude Agent SDK} to build agents that call Claude's API to make decisions:

\begin{infobox}
\textbf{Example:} Instead of a simple ``retry failed jobs'' script, a Claude-powered agent could:
\begin{itemize}[leftmargin=1.5em, itemsep=2pt]
  \item Analyze \textit{why} a generation failed (API timeout? content filter? malformed input?)
  \item Decide whether to retry, modify the prompt and retry, or flag for human review
  \item Draft a customer apology email if the delay exceeds 24 hours
  \item Adjust the story prompt to avoid whatever triggered the content filter
\end{itemize}
This is the power of AI agents vs.\ simple automation scripts. Your Mac Mini M2 has more than enough power to run these locally.
\end{infobox}

\section{Remote Access to Your Mac Mini}

You will want to check on your agents when you are away from home.

\step{7} Enable remote access:
\begin{actionbox}
\textbf{Option A: SSH (free, terminal only)}\\
\textbf{System Settings} $\rightarrow$ \textbf{General} $\rightarrow$ \textbf{Sharing} $\rightarrow$ enable \textbf{Remote Login}.\\
From any other computer: \texttt{ssh yourusername@your-mac-ip}\\[4pt]
\textbf{Option B: Tailscale (free, works anywhere)}\\
Install Tailscale (\url{https://tailscale.com}, free for personal use) on both the Mac Mini and your phone/laptop. This creates a secure private network so you can SSH into your Mac Mini from anywhere in the world, even behind a router.\\[4pt]
\textbf{Option C: Screen Sharing (free, full GUI)}\\
\textbf{System Settings} $\rightarrow$ \textbf{General} $\rightarrow$ \textbf{Sharing} $\rightarrow$ enable \textbf{Screen Sharing}.\\
Use the built-in Screen Sharing app on any other Mac to see the full desktop.
\end{actionbox}

\section{Mac Mini vs.\ Make.com: Migration Path}

\begin{table}[H]
\centering
\begin{tabularx}{\textwidth}{lX}
\toprule
\textbf{Phase} & \textbf{What To Do} \\
\midrule
\textbf{Launch (Week 1--4)} & Use Make.com for critical workflows (payment $\rightarrow$ generation). Run Mac Mini agents for monitoring, health checks, and daily digests. \\
\textbf{Stable (Month 2--3)} & Migrate order processing and fulfillment tracking to Mac Mini agents. Keep Make.com as a backup. \\
\textbf{Confident (Month 4+)} & Cancel Make.com entirely (\textbf{saves \$11--29/month}). All automation runs on your Mac Mini. \\
\bottomrule
\end{tabularx}
\end{table}

\begin{warningbox}
\textbf{Downside of self-hosted agents:} If your Mac Mini loses power or internet, your agents stop. Mitigate this with:
\begin{itemize}[leftmargin=1.5em]
  \item A UPS (uninterruptible power supply, \$40--80) for short outages
  \item The health monitor agent emailing you if it detects issues
  \item Keeping Make.com active for the first few months as a fallback
  \item Your backend on Railway still processes webhooks regardless---the Mac Mini agents are supplementary
\end{itemize}
\end{warningbox}

\section{Chapter 8 Cost Summary}

\begin{table}[H]
\centering
\begin{tabular}{lr}
\toprule
\textbf{Item} & \textbf{Cost} \\
\midrule
Mac Mini M2 (already owned) & \$0.00 \\
Electricity (\~{}10W idle, 24/7) & \$1--2/month \\
Tailscale (free personal tier) & \$0.00 \\
UPS battery backup (optional) & \$40--80 one-time \\
pm2 (process manager) & \$0.00 \\
\midrule
\textbf{Chapter 8 Total} & \textbf{\$0--80 one-time, \$1--2/month} \\
\bottomrule
\end{tabular}
\end{table}


% ══════════════════════════════════════════════════════════════
% CHAPTER 9: PRINT-ON-DEMAND
% ══════════════════════════════════════════════════════════════
\chapter{Connect Print-on-Demand (Lulu Direct)}

\section{How Print-on-Demand Works}

You never hold inventory. When a customer orders a book:

\begin{enumerate}[leftmargin=2em]
  \item Your system sends the print-ready PDF to Lulu's API with the customer's selected book options
  \item Lulu prints a \textbf{single copy} at their nearest facility using the specified binding, paper, and size
  \item Lulu ships it directly to your customer in \textbf{your brand's packaging} (white-label)
  \item You pay Lulu only for the printing + shipping of that one book
\end{enumerate}

\noindent There is \textbf{no minimum order} and \textbf{no upfront inventory cost}.

\section{Customer-Selectable Book Options}

Your book creation wizard (already built in \texttt{frontend/src/components/wizard/StepBookOptions.tsx}) lets customers choose from the following options. The price updates live in the browser as they customize:

\begin{table}[H]
\centering
\small
\begin{tabularx}{\textwidth}{lXrl}
\toprule
\textbf{Option} & \textbf{Choices} & \textbf{Surcharge} & \textbf{Default} \\
\midrule
\textbf{Page Count} & 12, 24, or 36 pages & Base price varies & 24 \\
\textbf{Binding} & Softcover (perfect-bound) & --- & Softcover \\
& Hardcover (casewrap) & +\$10.00 & \\
& Saddle Stitch (stapled, $\leq$32 pages) & --- & \\
\textbf{Paper} & Standard White (60\#) & --- & Standard \\
& Premium Matte & +\$5.00 & \\
& Glossy & +\$8.00 & \\
\textbf{Book Size} & 8\,'' $\times$ 8\,'' Square & --- & Square \\
& 8.5\,'' $\times$ 11\,'' Portrait & +\$3.00 & \\
& 11\,'' $\times$ 8.5\,'' Landscape & +\$3.00 & \\
\textbf{Illustration} & Full Color (AI-generated) & --- & Full Color \\
& Black \& White (coloring book) & --- & \\
\textbf{Add-ons} & Digital PDF download & +\$4.99 & Off \\
& Audiobook (AI-narrated) & +\$9.99 & Off \\
& Gift Wrap & +\$3.99 & Off \\
\bottomrule
\end{tabularx}
\end{table}

\begin{infobox}
\textbf{How pricing works:} The base price depends on page count: \$19.99 (12 pages), \$24.99 (24 pages), or \$34.99 (36 pages). Surcharges for binding, paper, size, and add-ons are stacked on top. Your backend (\texttt{books.service.ts}) calculates the exact price, and the frontend wizard shows it live as the customer makes selections.
\end{infobox}

\section{Compatibility Rules}

Not all combinations are physically possible. Your codebase enforces these rules in both the frontend (warning banners) and backend (validation errors):

\begin{warningbox}
\textbf{Incompatible combinations your system prevents:}
\begin{itemize}[leftmargin=1.5em, itemsep=2pt]
  \item \textbf{Saddle stitch + more than 32 pages} --- Too many pages for stapled binding. Customer must choose softcover/hardcover or reduce pages.
  \item \textbf{Softcover or hardcover + fewer than 24 pages} --- Perfect binding requires a minimum spine width. Customer must choose saddle stitch for 12-page books.
  \item \textbf{Hardcover + landscape format} --- Not available through Lulu Direct. Customer must choose portrait/square or softcover.
\end{itemize}
\end{warningbox}

\section{How Options Map to Lulu's API}

Your backend (\texttt{backend/src/fulfillment/providers/lulu.provider.ts}) automatically maps customer selections to Lulu's \texttt{pod\_package\_id} format when submitting a print job:

\begin{table}[H]
\centering
\small
\begin{tabularx}{\textwidth}{lX}
\toprule
\textbf{Customer Selection} & \textbf{Lulu pod\_package\_id Component} \\
\midrule
8\,'' $\times$ 8\,'' Square & \texttt{0800X0800} \\
8.5\,'' $\times$ 11\,'' Portrait & \texttt{0850X1100} \\
11\,'' $\times$ 8.5\,'' Landscape & \texttt{1100X0850} \\
Softcover (perfect-bound) & \texttt{FC} \\
Hardcover (casewrap) & \texttt{CW} \\
Saddle Stitch & \texttt{SS} \\
Standard paper + matte cover & \texttt{STDM} \\
Premium paper + matte cover & \texttt{PREM} \\
Premium paper + glossy cover & \texttt{PREG} \\
\bottomrule
\end{tabularx}
\end{table}

\noindent Example: A customer ordering a \textbf{8$\times$8 softcover book on glossy paper} results in \texttt{pod\_package\_id = 0800X0800FCPREG}. This mapping is handled automatically by your code---you do not need to calculate it manually.

\section{PDF Specifications}

\begin{table}[H]
\centering
\begin{tabularx}{\textwidth}{lX}
\toprule
\textbf{Specification} & \textbf{Value} \\
\midrule
File Format & PDF, 300 DPI, RGB (Lulu converts to CMYK) \\
Bleed & 0.125\,'' (9pt) on all sides \\
Gutter (inner margin) & Varies by binding: 36pt softcover, 54pt hardcover, 24pt saddle stitch \\
Interior & Full color or B\&W (based on customer's illustration style selection) \\
Fonts & Embedded in PDF \\
\bottomrule
\end{tabularx}
\end{table}

\begin{infobox}
\textbf{Already handled:} Your \texttt{PdfAssemblyService} (\texttt{backend/src/ai/pdf/pdf-assembly.service.ts}) automatically adjusts bleed, gutter margins, and page dimensions based on the customer's binding type and book size selections. It uses binding-aware margins so hardcover books have wider gutters for the casewrap.
\end{infobox}

\section{Lulu Volume Discounts}

As your business grows, per-unit printing costs decrease:

\begin{table}[H]
\centering
\begin{tabular}{lr}
\toprule
\textbf{Quantity Per Order} & \textbf{Discount} \\
\midrule
1--29 copies & 0\% (base price) \\
30--59 copies & 5\% off \\
60--99 copies & 10\% off \\
100--499 copies & 15\% off \\
500+ copies & 20\% off \\
\bottomrule
\end{tabular}
\end{table}

\section{Webhook Setup for Automatic Status Updates}

When Lulu ships a book, they can notify your system automatically via a ``webhook'' (an automatic notification).

\begin{actionbox}
In your Lulu developer dashboard, set your webhook URL to:\\[4pt]
\texttt{https://api.yourdomain.com/api/v1/fulfillment/webhook/lulu}\\[4pt]
Select events: \texttt{PRINT\_JOB.STATUS\_CHANGED}, \texttt{PRINT\_JOB.SHIPPED}\\[4pt]
Your backend already has the handler for this built in.
\end{actionbox}


% ══════════════════════════════════════════════════════════════
% CHAPTER 9: PAYMENTS & SUBSCRIPTIONS
% ══════════════════════════════════════════════════════════════
\chapter{Configure Payments and Subscriptions}

\section{Stripe Fee Structure}

Every transaction incurs fees. Build these into your pricing:

\begin{table}[H]
\centering
\begin{tabularx}{\textwidth}{Xr}
\toprule
\textbf{Transaction Type} & \textbf{Fee} \\
\midrule
Domestic credit/debit card & 2.9\% + \$0.30 \\
International card & 3.9\% + \$0.30 \\
Subscription billing surcharge & Additional 0.7\% \\
ACH bank transfer & 0.8\% (max \$5.00) \\
Refund & Original fee is NOT returned \\
Chargeback (disputed payment) & \$15.00 per dispute \\
\bottomrule
\end{tabularx}
\end{table}

\section{Example Fee Calculation}

For a \$29.99 book purchased with a domestic credit card:

\begin{table}[H]
\centering
\begin{tabular}{lr}
\toprule
Sale price & \$29.99 \\
Stripe fee (2.9\% + \$0.30) & --\$1.17 \\
\textbf{You receive} & \textbf{\$28.82} \\
\bottomrule
\end{tabular}
\end{table}

For a \$24.99/year subscription renewal:
\begin{table}[H]
\centering
\begin{tabular}{lr}
\toprule
Subscription price & \$24.99 \\
Stripe payment fee (2.9\% + \$0.30) & --\$1.02 \\
Stripe billing fee (0.7\%) & --\$0.17 \\
\textbf{You receive} & \textbf{\$23.80} \\
\bottomrule
\end{tabular}
\end{table}

\section{Setting Up Stripe Webhooks}

Webhooks let Stripe notify your app when payments succeed, fail, or get refunded.

\begin{actionbox}
In Stripe Dashboard $\rightarrow$ \textbf{Developers} $\rightarrow$ \textbf{Webhooks} $\rightarrow$ \textbf{Add endpoint}:\\[4pt]
URL: \texttt{https://api.yourdomain.com/api/v1/payments/webhook}\\[4pt]
Select these events:
\begin{itemize}[leftmargin=1.5em]
  \item \texttt{payment\_intent.succeeded}
  \item \texttt{payment\_intent.payment\_failed}
  \item \texttt{invoice.payment\_succeeded}
  \item \texttt{invoice.payment\_failed}
  \item \texttt{customer.subscription.deleted}
\end{itemize}
Copy the \textbf{Webhook Signing Secret} (starts with \texttt{whsec\_}) and add it to your backend environment variables as \texttt{STRIPE\_WEBHOOK\_SECRET}.
\end{actionbox}


% ══════════════════════════════════════════════════════════════
% CHAPTER 10: PER-BOOK COST ANALYSIS
% ══════════════════════════════════════════════════════════════
\chapter{Per-Book Cost Analysis}

This is the most important chapter for your business planning. Every number here is based on real pricing as of February 2026.

\section{AI Story Generation Costs: Claude vs.\ OpenAI}

Choosing the right AI model for story generation is critical---it directly affects the quality of your product. Children's book stories require creative prose, age-appropriate vocabulary, narrative coherence, and natural-sounding personalization. This section compares every viable API model as of February 2026.

\subsection{Token Usage for a 30-Page Book}

Each book generation requires several API calls. Here is the estimated token usage:

\begin{table}[H]
\centering
\small
\begin{tabularx}{\textwidth}{Xrrr}
\toprule
\textbf{Pipeline Stage} & \textbf{Input Tokens} & \textbf{Output Tokens} & \textbf{API Calls} \\
\midrule
Story outline generation & 800 & 1{,}200 & 1 \\
Page-by-page story text (30 pages) & 8{,}000 & 4{,}000 & 2--3 \\
Illustration prompt generation (16 images) & 4{,}900 & 2{,}400 & 1--2 \\
Content safety validation & 4{,}500 & 800 & 1 \\
Character description (from photo, optional) & 800 & 500 & 1 \\
\midrule
\textbf{Total} & \textbf{19{,}000} & \textbf{8{,}900} & \textbf{6--8} \\
\bottomrule
\end{tabularx}
\caption{Token usage for a 30-page personalized children's book}
\end{table}

\subsection{Cost Per Model: Full Comparison}

\begin{table}[H]
\centering
\small
\begin{tabularx}{\textwidth}{Xrrrr}
\toprule
\textbf{Model} & \textbf{Input \$/1M} & \textbf{Output \$/1M} & \textbf{Cost/Book} & \textbf{Creative Quality} \\
\midrule
\multicolumn{5}{l}{\textit{Anthropic Claude (recommended for story text):}} \\
\textbf{Claude Opus 4.6} & \$5.00 & \$25.00 & \cost{\$0.32} & Best in class \\
\textbf{Claude Sonnet 4.5} & \$3.00 & \$15.00 & \cost{\$0.19} & Excellent \\
Claude Haiku 4.5 & \$1.00 & \$5.00 & \cost{\$0.064} & Good \\
\midrule
\multicolumn{5}{l}{\textit{OpenAI GPT (good alternative for structured tasks):}} \\
GPT-4.1 & \$2.00 & \$8.00 & \cost{\$0.11} & Very good \\
GPT-4o & \$2.50 & \$10.00 & \cost{\$0.13} & Very good \\
GPT-4.1 Mini & \$0.40 & \$1.60 & \cost{\$0.022} & Acceptable \\
GPT-4o Mini & \$0.15 & \$0.60 & \cost{\$0.008} & Fair \\
GPT-4.1 Nano & \$0.10 & \$0.40 & \cost{\$0.006} & Low \\
\midrule
\multicolumn{5}{l}{\textit{OpenAI Reasoning Models (not recommended for creative writing):}} \\
o3 & \$2.00 & \$8.00 & \cost{\$0.15--0.30*} & Overkill \\
o4-mini & \$1.10 & \$4.40 & \cost{\$0.10--0.20*} & Overkill \\
\bottomrule
\end{tabularx}
\caption{Story generation cost per 30-page book by model. *O-series costs are unpredictable due to hidden reasoning tokens.}
\end{table}

\begin{warningbox}
\textbf{Avoid O-series reasoning models} (o1, o3, o4-mini) for story generation. These models consume invisible ``reasoning tokens'' billed as output that can 3--5$\times$ the actual cost. They are designed for math/logic problems, not creative writing. Use Claude or GPT-4.1 instead.
\end{warningbox}

\subsection{Why Claude Excels at Children's Stories}

In independent benchmarks and community evaluations, Claude models consistently outperform GPT models for creative fiction writing. Specifically:

\begin{itemize}[leftmargin=2em, itemsep=2pt]
  \item \textbf{Natural prose:} Claude produces more vivid, flowing narrative text that sounds like a real children's book author
  \item \textbf{Age-appropriate voice:} Claude better adapts vocabulary and sentence complexity to the target age range (e.g., simpler words for ages 3--5, richer language for ages 8--10)
  \item \textbf{Character consistency:} Claude maintains the child's personality traits and interests throughout the story without forced references
  \item \textbf{Moral lesson integration:} Claude weaves themes naturally into the narrative rather than adding preachy lessons at the end
  \item \textbf{Read-aloud quality:} Claude's prose has better rhythm and cadence---important for books parents read aloud
  \item \textbf{Personalization:} Claude integrates the child's name, interests, and details in a way that feels organic, not template-like
\end{itemize}

\begin{successbox}
\textbf{Recommended configuration:} Use \textbf{Claude Sonnet 4.5} for story outline and page-by-page text generation (where creative quality is paramount), and \textbf{GPT-4.1} for illustration prompt generation and safety validation (where structured output matters more than prose quality). This hybrid approach costs approximately \textbf{\$0.21 per book}---and with batch processing, as low as \textbf{\$0.11 per book}.
\end{successbox}

\subsection{Recommended Hybrid Pipeline}

\begin{table}[H]
\centering
\begin{tabularx}{\textwidth}{lXrr}
\toprule
\textbf{Stage} & \textbf{Model} & \textbf{Cost} & \textbf{Why This Model} \\
\midrule
Story outline & Claude Sonnet 4.5 & \$0.02 & Best narrative structure \\
Page-by-page text & Claude Sonnet 4.5 & \$0.15 & Best creative prose \\
Illustration prompts & GPT-4.1 & \$0.03 & Structured visual descriptions \\
Safety validation & GPT-4.1 Nano & \$0.003 & Simple classification, cheapest \\
Character description & GPT-4o (vision) & \$0.01 & Vision model for photo analysis \\
\midrule
\textbf{Total per book} & & \textbf{\$0.21} & \\
\textbf{With batch API (50\% off)} & & \textbf{\$0.11} & \\
\bottomrule
\end{tabularx}
\caption{Recommended hybrid pipeline: Claude for story quality, GPT for structured tasks}
\end{table}

\begin{infobox}
\textbf{Batch processing:} Both Anthropic and OpenAI offer a \textbf{Batch API} with 50\% off all token costs. Results are returned within 24 hours instead of immediately. Since book generation already takes 5--15 minutes (mostly waiting for illustrations), the batch API is a natural fit. Configure your backend to use batch endpoints for non-urgent generations.
\end{infobox}

\subsection{Alternative Configurations}

\begin{table}[H]
\centering
\small
\begin{tabularx}{\textwidth}{lXrr}
\toprule
\textbf{Config} & \textbf{Models Used} & \textbf{Cost/Book} & \textbf{Quality} \\
\midrule
\textbf{Premium} & Claude Opus 4.6 (all text) + GPT-4o (prompts) & \$0.37 & Best possible \\
\textbf{Recommended} & Claude Sonnet 4.5 (text) + GPT-4.1 (prompts) & \$0.21 & Excellent \\
\textbf{Budget} & Claude Haiku 4.5 (text) + GPT-4.1 Mini (prompts) & \$0.08 & Good \\
\textbf{Minimum} & GPT-4.1 (text) + GPT-4.1 Nano (prompts) & \$0.05 & Acceptable \\
\bottomrule
\end{tabularx}
\caption{Four configuration options from premium to minimum cost}
\end{table}

\begin{warningbox}
\textbf{Do not use the ``Minimum'' configuration for paying customers.} GPT-4.1 Nano produces noticeably lower-quality prose. It is acceptable for testing and development but not for a product parents are paying \$25--45 for. The difference between ``Recommended'' (\$0.21) and ``Minimum'' (\$0.05) is only \textbf{16 cents per book}---that is not worth sacrificing quality.
\end{warningbox}

\subsection{Anthropic API Note: Using Claude Without Extra Cost}

You already pay for \textbf{Claude Max}, which gives you unlimited Claude Code usage. However, Claude Max does \textbf{not} include API access for your backend. To use Claude Sonnet 4.5 in your production pipeline, you need a separate \textbf{Anthropic API account} at \url{https://console.anthropic.com}:

\begin{actionbox}
\step{1} Go to \url{https://console.anthropic.com} and create an account (or log in with your existing Anthropic account).\\[4pt]
\step{2} Add a payment method (your business debit card).\\[4pt]
\step{3} Add \cost{\$20} in prepaid credits. At \$0.21/book, this covers approximately \textbf{95 books} of story generation.\\[4pt]
\step{4} Create an API key. Save it in Bitwarden and add it to your \texttt{backend/.env} as \texttt{ANTHROPIC\_API\_KEY}.\\[4pt]
\step{5} Set a monthly spending limit of \cost{\$50} under Settings $\rightarrow$ Limits to protect yourself.
\end{actionbox}

\subsection{Image Generation (Illustrations)}

Every other page gets an illustration, plus a cover image:

\noindent \textbf{Image count by page count:}
\begin{itemize}[leftmargin=2em]
  \item 12-page book: 6 interior illustrations + 1 cover = \textbf{7 images}
  \item 24-page book: 12 interior illustrations + 1 cover = \textbf{13 images}
  \item 30-page book: 15 interior illustrations + 1 cover = \textbf{16 images}
  \item 36-page book: 18 interior illustrations + 1 cover = \textbf{19 images}
\end{itemize}

\begin{table}[H]
\centering
\small
\begin{tabularx}{\textwidth}{Xrrrrr}
\toprule
\textbf{Provider} & \textbf{Per Image} & \textbf{13 imgs} & \textbf{16 imgs} & \textbf{19 imgs} & \textbf{Quality} \\
& & \textbf{(24pg)} & \textbf{(30pg)} & \textbf{(36pg)} \\
\midrule
DALL-E 3 HD & \$0.080 & \$1.04 & \$1.28 & \$1.52 & Excellent \\
DALL-E 3 Standard & \$0.040 & \$0.52 & \$0.64 & \$0.76 & Very Good \\
Flux.2 Pro (BFL) & \$0.030 & \$0.39 & \$0.48 & \$0.57 & Excellent \\
Stability AI Core & \$0.030 & \$0.39 & \$0.48 & \$0.57 & Good \\
Leonardo AI & $\sim$\$0.015 & \$0.20 & \$0.24 & \$0.29 & Very Good \\
GPT Image Mini & \$0.005 & \$0.07 & \$0.08 & \$0.10 & Good \\
Flux.1 Schnell & \$0.003 & \$0.04 & \$0.05 & \$0.06 & Good \\
\bottomrule
\end{tabularx}
\caption{Image generation costs per book (every other page + cover)}
\end{table}

\section{Printing and Fulfillment Costs}

Printing costs vary based on the customer's binding, paper, and size selections. Here are the ranges for the default 8\,'' $\times$ 8\,'' format:

\begin{table}[H]
\centering
\begin{tabular}{lrr}
\toprule
\textbf{Component} & \textbf{24-Page Book} & \textbf{36-Page Book} \\
\midrule
Lulu print cost (softcover, standard paper) & \$3.00--5.00 & \$5.50--9.00 \\
Lulu print cost (hardcover, standard paper) & \$8.00--11.00 & \$10.50--14.00 \\
Lulu fulfillment fee & \$1.75 & \$1.75 \\
Shipping (standard domestic) & \$3.00--5.00 & \$3.50--6.00 \\
\midrule
\textbf{Softcover printing subtotal} & \textbf{\$7.75--11.75} & \textbf{\$10.75--16.75} \\
\textbf{Hardcover printing subtotal} & \textbf{\$12.75--17.75} & \textbf{\$15.75--21.75} \\
\bottomrule
\end{tabular}
\end{table}

\begin{infobox}
\textbf{Why the surcharges are profitable:} Hardcover adds \$10 to the retail price but only costs you \$5--6 more to print. Premium matte paper adds \$5 to retail but costs $\sim$\$1--2 more. \textbf{Upgrades increase your margins, not just your revenue.} Offering more options lets customers self-select into higher price tiers.
\end{infobox}

\section{Complete Cost-of-Goods-Sold (COGS) Per Book}

Here is the \textbf{total cost to produce and ship one book}, under three scenarios:

\subsection{Scenario A: Premium Quality (DALL-E 3 HD)}

\begin{table}[H]
\centering
\begin{tabular}{lrr}
\toprule
\textbf{Cost Component} & \textbf{24-Page} & \textbf{50-Page} \\
\midrule
AI text generation (GPT-4o) & \$0.10 & \$0.18 \\
AI images (DALL-E 3 HD, 13/26 imgs) & \$1.04 & \$2.08 \\
Lulu printing & \$4.00 & \$7.00 \\
Lulu fulfillment fee & \$1.75 & \$1.75 \\
Shipping (standard) & \$4.00 & \$4.50 \\
Stripe fee (on \$29.99 / \$44.99 sale) & \$1.17 & \$1.60 \\
\midrule
\textbf{Total COGS} & \textbf{\$12.06} & \textbf{\$17.11} \\
\textbf{Revenue at listed price} & \$29.99 & \$44.99 \\
\textbf{Gross Profit} & \textcolor{success}{\textbf{\$17.93}} & \textcolor{success}{\textbf{\$27.88}} \\
\textbf{Gross Margin} & \textcolor{success}{\textbf{59.8\%}} & \textcolor{success}{\textbf{62.0\%}} \\
\bottomrule
\end{tabular}
\caption{Scenario A: DALL-E 3 HD --- Premium quality}
\end{table}

\subsection{Scenario B: Balanced (Flux.2 Pro --- Recommended)}

\begin{table}[H]
\centering
\begin{tabular}{lrr}
\toprule
\textbf{Cost Component} & \textbf{24-Page} & \textbf{50-Page} \\
\midrule
AI text generation (GPT-4o) & \$0.10 & \$0.18 \\
AI images (Flux.2 Pro, 13/26 imgs) & \$0.39 & \$0.78 \\
Lulu printing & \$4.00 & \$7.00 \\
Lulu fulfillment fee & \$1.75 & \$1.75 \\
Shipping (standard) & \$4.00 & \$4.50 \\
Stripe fee (on \$29.99 / \$44.99 sale) & \$1.17 & \$1.60 \\
\midrule
\textbf{Total COGS} & \textbf{\$11.41} & \textbf{\$15.81} \\
\textbf{Revenue at listed price} & \$29.99 & \$44.99 \\
\textbf{Gross Profit} & \textcolor{success}{\textbf{\$18.58}} & \textcolor{success}{\textbf{\$29.18}} \\
\textbf{Gross Margin} & \textcolor{success}{\textbf{62.0\%}} & \textcolor{success}{\textbf{64.9\%}} \\
\bottomrule
\end{tabular}
\caption{Scenario B: Flux.2 Pro --- Recommended balance of quality and cost}
\end{table}

\subsection{Scenario C: Budget (Flux.1 Schnell + GPT-4o Mini)}

\begin{table}[H]
\centering
\begin{tabular}{lrr}
\toprule
\textbf{Cost Component} & \textbf{24-Page} & \textbf{50-Page} \\
\midrule
AI text generation (GPT-4o Mini) & \$0.01 & \$0.01 \\
AI images (Flux.1 Schnell, 13/26 imgs) & \$0.04 & \$0.08 \\
Lulu printing & \$4.00 & \$7.00 \\
Lulu fulfillment fee & \$1.75 & \$1.75 \\
Shipping (standard) & \$4.00 & \$4.50 \\
Stripe fee (on \$29.99 / \$44.99 sale) & \$1.17 & \$1.60 \\
\midrule
\textbf{Total COGS} & \textbf{\$10.97} & \textbf{\$14.94} \\
\textbf{Revenue at listed price} & \$29.99 & \$44.99 \\
\textbf{Gross Profit} & \textcolor{success}{\textbf{\$19.02}} & \textcolor{success}{\textbf{\$30.05}} \\
\textbf{Gross Margin} & \textcolor{success}{\textbf{63.4\%}} & \textcolor{success}{\textbf{66.8\%}} \\
\bottomrule
\end{tabular}
\caption{Scenario C: Budget --- Lowest possible AI costs}
\end{table}

\begin{successbox}
\textbf{Key Insight:} AI generation costs are almost negligible compared to printing and shipping. The difference between the cheapest and most expensive AI option is only \$1.00--2.00 per book. \textbf{Do not sacrifice quality to save \$1.} Use at least Flux.2 Pro or DALL-E 3 Standard for professional-looking illustrations.
\end{successbox}


% ══════════════════════════════════════════════════════════════
% CHAPTER 11: STARTUP BUDGET
% ══════════════════════════════════════════════════════════════
\chapter{Complete Startup Budget Breakdown}

Your total budget is \$2{,}000. Here is exactly how to allocate it.

\section{One-Time Startup Costs}

\begin{table}[H]
\centering
\begin{tabularx}{\textwidth}{Xlr}
\toprule
\textbf{Expense} & \textbf{Category} & \textbf{Amount} \\
\midrule
SC LLC formation (online filing) & Legal & \$125.00 \\
Domain name (1 year via Cloudflare) & Infrastructure & \$10.44 \\
OpenAI API prepaid credits & AI Services & \$20.00 \\
Test book printing (3 samples from Lulu) & Testing & \$45.00 \\
Initial marketing (social media ads) & Marketing & \$200.00 \\
Logo design (Canva Pro 1 month, optional) & Brand & \$13.00 \\
\midrule
\textbf{One-Time Total} & & \textbf{\$413.44} \\
\bottomrule
\end{tabularx}
\end{table}

\section{Monthly Operating Costs (Months 1--3, Pre-Revenue)}

\begin{table}[H]
\centering
\begin{tabularx}{\textwidth}{Xlr}
\toprule
\textbf{Expense} & \textbf{Category} & \textbf{Monthly} \\
\midrule
Railway.app hosting (Hobby) & Infrastructure & \$5.00 \\
Railway compute/database overage & Infrastructure & \$5.00--15.00 \\
Make.com (Core plan) & Automation & \$10.59 \\
Claude Code (via Claude Max) & Development & \$0.00 (already subscribed) \\
SendGrid (free tier) & Email & \$0.00 \\
OpenAI API (testing/generation) & AI Services & \$5.00--10.00 \\
\midrule
\textbf{Monthly Total (pre-revenue)} & & \textbf{\$25.59--35.59} \\
\bottomrule
\end{tabularx}
\end{table}

\section{Budget Runway Calculation}

\begin{table}[H]
\centering
\begin{tabular}{lr}
\toprule
Total budget & \$2{,}000.00 \\
Minus one-time costs & --\$413.44 \\
\midrule
Remaining for operations & \$1{,}586.56 \\
Monthly burn rate (estimated) & \$30.00 \\
\midrule
\textbf{Months of runway (zero revenue)} & \textbf{52.9 months} \\
\bottomrule
\end{tabular}
\end{table}

\begin{successbox}
\textbf{You have over 4 years of runway} before you need a single sale. Because your Claude Max subscription eliminates AI coding tool costs, your monthly burn is remarkably low. Take your time to get the product right, test thoroughly, and build an audience before revenue pressure hits.
\end{successbox}

\section{When Revenue Starts: Variable Costs Per Sale}

Once you start selling, each sale generates revenue \textit{and} costs. The customer's selections affect both the retail price and your cost. Here are two common configurations:

\begin{table}[H]
\centering
\small
\begin{tabular}{lrr}
\toprule
& \textbf{24pg Softcover Std} & \textbf{24pg Hardcover Premium} \\
& \textbf{(\$24.99 retail)} & \textbf{(\$39.99 retail)} \\
\midrule
Revenue per sale & \$24.99 & \$39.99 \\
AI generation (Flux.2 Pro) & --\$0.49 & --\$0.49 \\
Printing (Lulu, varies by binding/paper) & --\$4.00 & --\$9.00 \\
Fulfillment fee (Lulu) & --\$1.75 & --\$1.75 \\
Shipping & --\$4.00 & --\$4.00 \\
Stripe fee & --\$1.02 & --\$1.46 \\
\midrule
\textbf{Net profit per sale} & \textcolor{success}{\textbf{\$13.73}} & \textcolor{success}{\textbf{\$23.29}} \\
\textbf{Gross margin} & \textcolor{success}{\textbf{54.9\%}} & \textcolor{success}{\textbf{58.2\%}} \\
\bottomrule
\end{tabular}
\end{table}

\begin{infobox}
\textbf{Upgrade economics:} When a customer upgrades from the base softcover (\$24.99) to a hardcover with premium matte paper (\$39.99), your profit jumps from \$13.73 to \$23.29---a \textbf{70\% increase in profit} per sale. This is why offering customization options is so valuable. Let customers self-select into premium tiers.
\end{infobox}

\section{Break-Even Analysis}

How many books must you sell to cover your monthly fixed costs?

\begin{table}[H]
\centering
\begin{tabular}{lr}
\toprule
Monthly fixed costs & \$30.00 \\
Profit per 24-page book sale & \$18.58 \\
\midrule
\textbf{Break-even: books per month} & \textbf{1.6 $\approx$ 2 books} \\
\bottomrule
\end{tabular}
\end{table}

\begin{successbox}
\textbf{You break even at just 2 sales per month.} That is extremely achievable. At 10 sales/month, you net approximately \$156/month in profit after all costs. At 50 sales/month, approximately \$900/month.
\end{successbox}


% ══════════════════════════════════════════════════════════════
% CHAPTER 12: MONTHLY OPERATING COSTS AT SCALE
% ══════════════════════════════════════════════════════════════
\chapter{Monthly Operating Costs at Scale}

\begin{table}[H]
\centering
\small
\begin{tabularx}{\textwidth}{Xrrrr}
\toprule
\textbf{Expense} & \textbf{0 orders} & \textbf{50/mo} & \textbf{200/mo} & \textbf{500/mo} \\
\midrule
Railway hosting & \$10 & \$20 & \$40 & \$80 \\
Make.com & \$11 & \$11 & \$19 & \$29 \\
Claude Code (Claude Max) & \$0 & \$0 & \$0 & \$0 \\
SendGrid & \$0 & \$0 & \$20 & \$20 \\
OpenAI API & \$0 & \$25 & \$98 & \$245 \\
Domain renewal & \$1 & \$1 & \$1 & \$1 \\
\midrule
\textbf{Fixed costs} & \textbf{\$22} & \textbf{\$57} & \textbf{\$178} & \textbf{\$375} \\
\midrule
Printing (50 $\times$ \$4) & --- & \$200 & \$800 & \$2{,}000 \\
Fulfillment fees & --- & \$88 & \$350 & \$875 \\
Shipping & --- & \$200 & \$800 & \$2{,}000 \\
Stripe fees & --- & \$59 & \$234 & \$585 \\
\midrule
\textbf{Total variable costs} & --- & \$547 & \$2{,}184 & \$5{,}460 \\
\midrule
\textbf{Total revenue} (at \$29.99) & \$0 & \$1{,}500 & \$5{,}998 & \$14{,}995 \\
\textbf{Total costs} & \$22 & \$604 & \$2{,}362 & \$5{,}835 \\
\textbf{Net Profit} & \textcolor{danger}{--\$22} & \textcolor{success}{\$896} & \textcolor{success}{\$3{,}636} & \textcolor{success}{\$9{,}160} \\
\textbf{Profit Margin} & --- & 59.7\% & 60.6\% & 61.1\% \\
\bottomrule
\end{tabularx}
\caption{Monthly P\&L at different order volumes (24-page books at \$29.99)}
\end{table}


% ══════════════════════════════════════════════════════════════
% CHAPTER 13: PRICING STRATEGY
% ══════════════════════════════════════════════════════════════
\chapter{Pricing Strategy}

\section{Recommended Retail Prices}

Your pricing is \textbf{dynamic}---the frontend wizard calculates the price live as customers select options. The backend verifies the price before creating the Stripe payment intent. Here are the base prices and upgrade surcharges:

\begin{table}[H]
\centering
\begin{tabularx}{\textwidth}{Xrrr}
\toprule
\textbf{Product / Option} & \textbf{Price} & \textbf{COGS} & \textbf{Margin} \\
\midrule
12-page softcover book (base) & \$19.99 & \$9.50 & 52\% \\
24-page softcover book (base) & \$24.99 & \$11.41 & 54\% \\
36-page softcover book (base) & \$34.99 & \$15.81 & 55\% \\
\midrule
\textit{Upgrade surcharges:} & & & \\
\quad Hardcover binding & +\$10.00 & +\$5.00 & 50\%+ \\
\quad Premium matte paper & +\$5.00 & +\$1.50 & 70\% \\
\quad Glossy paper \& cover & +\$8.00 & +\$2.00 & 75\% \\
\quad Larger size (8.5$\times$11 or landscape) & +\$3.00 & +\$1.00 & 67\% \\
\midrule
\textit{Add-ons:} & & & \\
\quad Digital PDF add-on & +\$4.99 & \$0.02 & 99\% \\
\quad Audiobook add-on & +\$9.99 & \$0.50 & 95\% \\
\quad Gift wrap & +\$3.99 & \$1.50 & 62\% \\
\midrule
Birthday subscription (yearly) & \$24.99/yr & \$11.41 & 54\% \\
Quarterly subscription & \$22.99/qtr & \$11.41 & 50\% \\
\bottomrule
\end{tabularx}
\end{table}

\begin{successbox}
\textbf{Price range example:} A basic 12-page saddle-stitch book on standard paper costs \$19.99. A fully loaded 36-page hardcover on glossy paper with audiobook, digital PDF, and gift wrap costs \$71.95. Most customers will land in the \$24.99--\$44.99 range.
\end{successbox}

\section{Upsell Strategy}

Your wizard design naturally encourages upgrades. The highest-margin items are displayed prominently:

\begin{itemize}[leftmargin=2em]
  \item \textbf{Hardcover upgrade} (+\$10.00): Your most profitable upsell. Costs you only \$5 more but adds \$10 to the price. Emphasize ``keepsake quality'' and ``lasts forever.''
  \item \textbf{Paper upgrades} (+\$5--8): Premium matte and glossy are high-margin. ``Feel the difference'' messaging works well.
  \item \textbf{Digital PDF} (+\$4.99): Costs you \$0.02 to generate. Almost pure profit. ``Read it on any device.''
  \item \textbf{Audiobook} (+\$9.99): Use OpenAI's text-to-speech API at approximately \$0.50/book. ``Let them hear their story.''
  \item \textbf{Gift wrap} (+\$3.99): Lulu may charge \$1--2 extra for packaging. ``Ready to give.''
  \item \textbf{Multiple copies}: Offer 10\% off for 2+ copies. Printing cost barely changes.
\end{itemize}

\section{Competitive Pricing Analysis}

\begin{table}[H]
\centering
\begin{tabularx}{\textwidth}{Xr}
\toprule
\textbf{Competitor} & \textbf{Price (comparable book)} \\
\midrule
Wonderbly (Lost My Name) & \$34.99--44.99 \\
Hooray Heroes & \$39.99--49.99 \\
I See Me! & \$34.99--39.99 \\
Shutterfly Photo Book & \$29.99--49.99 \\
\textbf{Crayons \& Quills (you)} & \textbf{\$19.99--71.95} \\
\bottomrule
\end{tabularx}
\end{table}

\noindent Your entry price (\$19.99) undercuts every competitor, while your premium configurations (\$35--72) compete with established players. Your AI-generated personalization offers a \textbf{unique value proposition} they cannot match (fully custom stories and illustrations, not template fill-in-the-blank). Your customization options let customers choose exactly the product they want---something most competitors do not offer.


% ══════════════════════════════════════════════════════════════
% CHAPTER 14: LAUNCH CHECKLIST
% ══════════════════════════════════════════════════════════════
\chapter{Launch Checklist}

Complete every item before announcing your business publicly.

\section{Pre-Launch Testing (Week of Launch)}

\begin{itemize}[leftmargin=2em, itemsep=4pt]
  \item[$\square$] \textbf{Order a test book yourself.} Go through the entire flow: create a child profile, configure a book, pay with Stripe (use test mode first, then a real \$1 charge), confirm the book generates, confirm the PDF looks correct, confirm Lulu receives the order, wait for the physical book to arrive. \textbf{This is the most important test.}
  \item[$\square$] Order 2--3 books with different options (different ages, themes, page counts)
  \item[$\square$] Test on a phone (not just a computer)---most customers will use mobile
  \item[$\square$] Test with a friend or family member who has never seen the site
  \item[$\square$] Verify all emails send correctly (order confirmation, shipping update)
  \item[$\square$] Verify Stripe webhook is receiving events (check Stripe Dashboard $\rightarrow$ Developers $\rightarrow$ Webhooks $\rightarrow$ Recent events)
  \item[$\square$] Verify Make.com automations trigger correctly (check execution history) and/or Mac Mini agents are running (\texttt{pm2 status})
  \item[$\square$] Test the refund flow (process a test refund through Stripe)
  \item[$\square$] Check all landing pages load correctly and have proper SEO meta tags
  \item[$\square$] Verify your domain has HTTPS (padlock icon in browser)
  \item[$\square$] Set up Google Analytics (free) at \url{https://analytics.google.com}
\end{itemize}

\section{Legal Compliance}

\begin{itemize}[leftmargin=2em, itemsep=4pt]
  \item[$\square$] Add a \textbf{Privacy Policy} page (use a free generator: \url{https://www.privacypolicygenerator.info})
  \item[$\square$] Add a \textbf{Terms of Service} page (free generator: \url{https://www.termsofservicegenerator.net})
  \item[$\square$] Add a \textbf{COPPA notice} explaining you collect children's names/ages only with parental consent and do not create accounts for children
  \item[$\square$] Add a \textbf{cookie consent} banner if using analytics
  \item[$\square$] Add a \textbf{refund policy} (recommendation: 30-day satisfaction guarantee)
\end{itemize}

\section{Go-Live Steps}

\begin{enumerate}[leftmargin=2em, itemsep=4pt]
  \item Switch Stripe from test mode to \textbf{live mode}
  \item Update all API keys in Railway to live keys
  \item Verify DNS is pointing correctly (domain resolves to your site)
  \item Announce on personal social media
  \item Post in relevant Facebook groups (parenting groups, gift idea groups)
  \item Submit to Product Hunt (\url{https://www.producthunt.com})
  \item Set up a simple Google Ads campaign (\$5--10/day budget)
\end{enumerate}


% ══════════════════════════════════════════════════════════════
% CHAPTER 15: ONGOING OPERATIONS
% ══════════════════════════════════════════════════════════════
\chapter{Ongoing Operations and Monitoring}

\section{Your Weekly Routine (2--5 Hours Total)}

\begin{table}[H]
\centering
\begin{tabularx}{\textwidth}{llX}
\toprule
\textbf{Day} & \textbf{Time} & \textbf{Task} \\
\midrule
Monday & 30 min & Check Stripe dashboard for revenue, failed payments, disputes \\
Monday & 15 min & Check Make.com / Mac Mini agent logs (\texttt{pm2 logs}) for failed automations \\
Monday & 15 min & Check Railway dashboard for any service errors \\
Wednesday & 30 min & Respond to any customer support emails \\
Wednesday & 30 min & Review any flagged/failed book generations \\
Friday & 15 min & Check Lulu dashboard for fulfillment status \\
Friday & 30 min & Marketing: schedule social media posts, adjust ad spend \\
Monthly & 1 hour & Review Wave accounting, categorize expenses, check P\&L \\
Quarterly & 2 hours & Pay estimated taxes, review with CPA \\
\bottomrule
\end{tabularx}
\end{table}

\section{What the AI Agents Handle Without You}

These processes run 24/7/365 with zero intervention:

\begin{enumerate}[leftmargin=2em, itemsep=2pt]
  \item Customer places order $\rightarrow$ Stripe processes payment
  \item Payment succeeds $\rightarrow$ Make.com or Mac Mini agent triggers book generation
  \item AI generates story $\rightarrow$ AI generates illustrations $\rightarrow$ PDF assembled
  \item PDF sent to Lulu $\rightarrow$ Book printed and shipped
  \item Customer receives confirmation email, then shipping email with tracking
  \item Subscription renewals auto-bill and auto-generate new books
  \item Abandoned carts get follow-up emails after 24 hours
  \item Fulfillment statuses sync every 4 hours
\end{enumerate}

\section{When You Need to Intervene}

Set up the ``Failed Generation Alert'' workflow (Chapter 7) so you get emailed when:
\begin{itemize}[leftmargin=2em]
  \item A book generation fails (usually an OpenAI API timeout---just retry)
  \item A payment fails (Stripe handles retries, but check for patterns)
  \item A fulfillment order is rejected by Lulu (usually a PDF formatting issue)
  \item A customer requests a refund (process via Stripe dashboard or admin API)
\end{itemize}

\section{Scaling Triggers}

When you hit these milestones, take the corresponding action:

\begin{table}[H]
\centering
\begin{tabularx}{\textwidth}{lX}
\toprule
\textbf{Milestone} & \textbf{Action} \\
\midrule
10 orders/month & Upgrade Railway to Pro (\$20/mo). You can afford it. \\
50 orders/month & Hire a virtual assistant (\$200--400/mo) for customer support \\
100 orders/month & Upgrade SendGrid to Essentials (\$20/mo) for higher email volume \\
200 orders/month & Consider adding a second print partner (Amazon KDP) for redundancy \\
500 orders/month & Hire a part-time developer for feature improvements \\
\bottomrule
\end{tabularx}
\end{table}


% ══════════════════════════════════════════════════════════════
% APPENDIX A: ACCOUNT TRACKER
% ══════════════════════════════════════════════════════════════
\appendix
\chapter{Complete Account Tracker}

Store all credentials in Bitwarden. This is your reference list:

\begin{longtable}{p{1.5in}p{1.8in}p{2.2in}}
\toprule
\textbf{Service} & \textbf{URL} & \textbf{Keys/Secrets Needed} \\
\midrule
\endhead
SC Secretary of State & businessfilings.sc.gov & Filing number \\
IRS & irs.gov & EIN \\
Mercury Bank & mercury.com & Account number, routing number \\
Bitwarden & bitwarden.com & Master password (memorize this) \\
Cloudflare & cloudflare.com & Account login \\
OpenAI & platform.openai.com & API Key (\texttt{sk-...}) \\
Stripe & stripe.com & Publishable Key (\texttt{pk\_live\_...}), Secret Key (\texttt{sk\_live\_...}), Webhook Secret (\texttt{whsec\_...}) \\
Lulu Direct & developers.lulu.com & Client Key, Client Secret \\
SendGrid & sendgrid.com & API Key (\texttt{SG....}) \\
Railway & railway.app & Login via GitHub \\
GitHub & github.com & Account login, SSH key \\
Make.com & make.com & Account login \\
Claude Code (Claude Max) & claude.ai/code & Already logged in via Claude Max \\
Wave Accounting & waveapps.com & Account login \\
Google Analytics & analytics.google.com & Measurement ID (\texttt{G-...}) \\
\bottomrule
\end{longtable}


% ══════════════════════════════════════════════════════════════
% APPENDIX B: AUTOMATION WORKFLOW DIAGRAMS
% ══════════════════════════════════════════════════════════════
\chapter{AI Agent Workflow Specifications}

\section{Master Automation Architecture}

\begin{center}
\begin{tikzpicture}[
  node distance=0.8cm and 1.5cm,
  box/.style={rectangle, draw=primary, fill=primary!10, text width=2.8cm, minimum height=1cm, align=center, rounded corners=3pt, font=\small},
  trigger/.style={rectangle, draw=accent, fill=accent!10, text width=2.5cm, minimum height=0.8cm, align=center, rounded corners=3pt, font=\small},
  arrow/.style={->, thick, color=midgray}
]

\node[trigger] (stripe) {Stripe Payment Event};
\node[box, right=of stripe] (generate) {Generate Book\\(OpenAI)};
\node[box, right=of generate] (lulu) {Submit to\\Lulu Print};
\node[box, right=of lulu] (email) {Email\\Customer};

\draw[arrow] (stripe) -- (generate);
\draw[arrow] (generate) -- (lulu);
\draw[arrow] (lulu) -- (email);

\node[trigger, below=1.5cm of stripe] (schedule) {Hourly\\Schedule};
\node[box, right=of schedule] (check) {Check Lulu\\Status};
\node[box, right=of check] (update) {Update\\Order DB};
\node[box, right=of update] (notify) {Ship\\Notification};

\draw[arrow] (schedule) -- (check);
\draw[arrow] (check) -- (update);
\draw[arrow] (update) -- (notify);

\node[trigger, below=1.5cm of schedule] (daily) {Daily\\9 AM};
\node[box, right=of daily] (abandoned) {Find Abandoned\\Carts};
\node[box, right=of abandoned] (recover) {Send Recovery\\Email};

\draw[arrow] (daily) -- (abandoned);
\draw[arrow] (abandoned) -- (recover);

\end{tikzpicture}
\end{center}

\section{Workflow Configuration Reference}

For each workflow in Make.com, here are the exact module configurations:

\subsection{Workflow 1: Order $\rightarrow$ Book $\rightarrow$ Print $\rightarrow$ Email}
\begin{enumerate}[leftmargin=2em]
  \item \textbf{Trigger:} Stripe $\rightarrow$ Watch Events $\rightarrow$ Event type: \texttt{payment\_intent.succeeded}
  \item \textbf{HTTP Module:} POST to \texttt{https://api.yourdomain.com/api/v1/books/\{\{bookId\}\}/generate}
  \begin{itemize}[leftmargin=1.5em]
    \item Header: \texttt{Authorization: Bearer \{\{adminToken\}\}}
  \end{itemize}
  \item \textbf{Sleep Module:} Wait 600 seconds (10 minutes) for generation
  \item \textbf{HTTP Module:} GET \texttt{https://api.yourdomain.com/api/v1/books/\{\{bookId\}\}}
  \item \textbf{Router:}
  \begin{itemize}[leftmargin=1.5em]
    \item Route 1 (status = REVIEW\_READY): Continue to fulfillment
    \item Route 2 (status = FAILED): Send alert email to you
  \end{itemize}
  \item \textbf{HTTP Module:} POST \texttt{https://api.yourdomain.com/api/v1/fulfillment/orders/\{\{orderId\}\}/submit}
  \item \textbf{SendGrid Module:} Send ``Order Confirmed'' template to customer email
\end{enumerate}


% ══════════════════════════════════════════════════════════════
% APPENDIX C: GLOSSARY
% ══════════════════════════════════════════════════════════════
\chapter{Glossary of Technical Terms}

If you encounter an unfamiliar term anywhere in this guide, look it up here.

\begin{description}[style=nextline, leftmargin=2em, itemsep=6pt]

\item[API (Application Programming Interface)]
A way for two software programs to talk to each other. When your website ``calls the OpenAI API,'' it is sending a request to OpenAI's servers and getting a response back. Think of it like ordering at a restaurant: you tell the waiter (API) what you want, the kitchen (server) makes it, and the waiter brings it back.

\item[API Key]
A secret password that identifies your account when using an API. Like a membership card that lets you into a club. Never share it publicly.

\item[Backend]
The part of your application that runs on a server (not visible to customers). It handles business logic, talks to the database, and communicates with external services like OpenAI and Stripe.

\item[BullMQ]
A software library that manages a queue of jobs. When a customer orders a book, the job ``generate this book'' goes into the queue. A worker process picks up jobs one by one and completes them.

\item[COGS (Cost of Goods Sold)]
The direct cost to produce one unit of your product. For you, this includes AI generation, printing, shipping, and payment processing fees.

\item[COPPA (Children's Online Privacy Protection Act)]
A US federal law that restricts what data websites can collect from children under 13. Since parents (not children) use your site, you are likely compliant, but you must clearly state this in your privacy policy.

\item[Deploy / Deployment]
Making your application available on the internet. ``Deploying to Railway'' means uploading your code to Railway's servers so anyone can visit your website.

\item[DNS (Domain Name System)]
The system that translates human-readable domain names (like \texttt{crayonsandquills.com}) into computer-readable IP addresses. You configure DNS records in Cloudflare to point your domain to your hosting provider.

\item[EIN (Employer Identification Number)]
A 9-digit number assigned by the IRS to identify your business for tax purposes. Like a Social Security number, but for your LLC.

\item[Environment Variable]
A configuration value stored outside your code. For example, your OpenAI API key is stored as an environment variable so it is not visible in the code itself. Think of it as a locked settings file.

\item[Frontend]
The part of your application that customers see and interact with in their web browser---the website itself.

\item[Git / GitHub]
Git is a tool that tracks changes to your code (like ``Track Changes'' in Microsoft Word). GitHub is a website that stores your code online and lets tools like Railway access it.

\item[Gross Margin]
The percentage of revenue remaining after subtracting COGS. A 60\% gross margin means you keep 60 cents of every dollar after production costs.

\item[HTTPS / SSL / TLS]
Encryption that protects data sent between a customer's browser and your server. Shown by the padlock icon in the browser address bar. Required for any site handling payments.

\item[JWT (JSON Web Token)]
A secure token (a long string of characters) that proves a user is logged in. When someone logs into your site, they receive a JWT that gets sent with every subsequent request.

\item[NestJS]
A framework (pre-built structure) for building backend applications in JavaScript/TypeScript. Your backend is built with NestJS.

\item[Next.js]
A framework for building frontend websites in React (JavaScript). Your frontend is built with Next.js.

\item[Node.js]
The engine that runs JavaScript code on a server (outside of a web browser). Both your backend and frontend require Node.js.

\item[PostgreSQL]
A database system that stores all your data (users, books, orders) in organized tables. Think of it as a very powerful spreadsheet.

\item[Print-on-Demand (POD)]
A business model where products are printed only after a customer orders them. No inventory, no upfront printing costs.

\item[Prisma]
A tool that makes it easy for your application to talk to the PostgreSQL database. Instead of writing complex database queries, you write simple commands.

\item[Redis]
An extremely fast, temporary data store. Your application uses it as a message queue: ``Here is a book that needs to be generated''---Redis holds that message until a worker picks it up.

\item[Repository (Repo)]
A folder of code stored on GitHub. Your entire Crayons \& Quills application lives in one repository.

\item[REST API]
A common way to structure an API using standard web addresses (URLs). For example, \texttt{GET /api/v1/books} retrieves all books, \texttt{POST /api/v1/books} creates a new one.

\item[Saddle-Stitch]
A binding method where pages are folded and stapled together. Common for thin booklets and children's books under 32 pages.

\item[SaaS (Software as a Service)]
A business model where customers pay to use software online (rather than buying a physical product). Your subscription plans make Crayons \& Quills a SaaS business.

\item[Stripe]
A payment processing company. They handle credit card charges, subscriptions, and payouts to your bank account.

\item[Token (AI)]
A unit of text processed by an AI model. Roughly 1 token = 4 characters or 0.75 words. AI APIs charge per token.

\item[Webhook]
An automatic notification sent from one service to another when something happens. When Stripe processes a payment, it sends a webhook to your backend saying ``payment succeeded.''

\item[Zustand]
A lightweight state management library for React. It stores the customer's selections as they go through the book creation wizard.

\end{description}

% ══════════════════════════════════════════════════════════════
% APPENDIX D: TIMELINE
% ══════════════════════════════════════════════════════════════
\chapter{Suggested 8-Week Launch Timeline}

\begin{table}[H]
\centering
\begin{tabularx}{\textwidth}{lX}
\toprule
\textbf{Week} & \textbf{Tasks} \\
\midrule
\textbf{Week 1} & Form LLC, get EIN, open business bank account, register domain \\
\textbf{Week 2} & Set up all technology accounts (Chapter 4), open Claude Code, clone your existing \texttt{LLM-books} codebase \\
\textbf{Week 3} & Configure \texttt{.env} files with your API keys, run Prisma migration, deploy to Railway, get site live \\
\textbf{Week 4} & Set up Stripe products and webhooks, configure Make.com workflows, set up Mac Mini agents \\
\textbf{Week 5} & Connect Lulu Direct, order 3 test books with different options, test full order flow end-to-end \\
\textbf{Week 6} & Fix any issues found in testing, add legal pages, set up analytics, refine agent scripts \\
\textbf{Week 7} & Soft launch: share with friends/family, gather feedback, iterate \\
\textbf{Week 8} & Public launch: social media, Product Hunt, initial ad campaigns \\
\bottomrule
\end{tabularx}
\end{table}

\begin{infobox}
\textbf{Do not rush Week 5.} Order real test books and wait for them to arrive. Hold the physical book in your hands. Check the print quality, colors, binding, and text readability. This is your product---it must be excellent before you sell it to anyone.
\end{infobox}

% ══════════════════════════════════════════════════════════════
% BACK MATTER
% ══════════════════════════════════════════════════════════════
\chapter*{Final Notes}
\addcontentsline{toc}{chapter}{Final Notes}

\begin{enumerate}[leftmargin=2em, itemsep=8pt]
  \item \textbf{Start simple.} Launch with 24-page bedtime stories for ages 4--7 only. Add more options after you have paying customers and real feedback.

  \item \textbf{Quality over speed.} A beautiful, well-printed 24-page book will generate referrals and repeat customers. A rushed, low-quality book will generate refunds and bad reviews.

  \item \textbf{The AI handles production. You handle marketing.} The bottleneck in this business is not technology---it is getting customers to your website. Invest your time in:
  \begin{itemize}[leftmargin=1.5em]
    \item Posting in parenting Facebook groups
    \item Creating TikTok/Instagram content showing the book creation process
    \item Running targeted Google Ads for ``personalized children's book''
    \item Building an email list with a free coloring page lead magnet
    \item Partnering with mommy bloggers and children's book reviewers
  \end{itemize}

  \item \textbf{Every dollar of your \$2{,}000 is accounted for.} You have 42+ months of runway at \$31/month. Use that time wisely---do not panic-spend on tools or services you do not yet need.

  \item \textbf{Reinvest profits.} Your first \$500 in revenue should go right back into advertising. Paid ads with a product that has 60\%+ margins create a flywheel: more ads $\rightarrow$ more sales $\rightarrow$ more profit $\rightarrow$ more ads.

  \item \textbf{Your codebase is your head start.} Most founders spend months building their product. Yours is already built. The frontend, backend, AI services, database schema, payment integration, and email system are all in your repository. Your job now is to configure, deploy, and market. That is a massive advantage.

  \item \textbf{You can do this.} The technology is ready. The codebase is built. The market is large. The margins are strong. Configure your accounts, deploy your site, order a test book, and launch. Execute the plan.
\end{enumerate}

\vfill
\begin{center}
\rule{2in}{0.4pt}\\[12pt]
{\large\textcolor{primary}{\textbf{Good luck. Go launch something wonderful.}}}
\end{center}

\end{document}
